\documentclass[11pt,a4paper]{article}

\usepackage[T2A]{fontenc}
\usepackage[utf8]{inputenc}
\usepackage[russian]{babel}
\usepackage[a4paper,margin=20mm]{geometry}
\usepackage{amsmath,amssymb,bm}
\usepackage{booktabs,tabularx,array}
\usepackage{microtype}
\usepackage[hidelinks]{hyperref}
\usepackage{mathrsfs}
\usepackage{tikz}
\usetikzlibrary{arrows.meta}

\newcommand{\PA}{P_A}
\newcommand{\VA}{V_A}
\newcommand{\Rref}{\mathcal R}
\newcommand{\dd}{\mathrm d}
\newcommand{\ph}{\mathrm{ph}}
\providecommand{\U}{}
\renewcommand{\U}{\bm U}
\newcommand{\F}{\bm F}
\newcommand{\Src}{\bm S}
\setlength{\parindent}{0pt}
\setlength{\parskip}{0.4em}
\setlength{\tabcolsep}{5pt}
\renewcommand{\arraystretch}{1.08}
\setlength{\emergencystretch}{2em}

\begin{document}

\begin{center}
{\Large\bfseries Одномерная нестационарная модель солнечного ветра}\\[4pt]
\end{center}

\section{Исходная постановка}

Рассматриваем течение вдоль открытой магнитной трубки
с заданной площадью $A(r)>0$,
$R_\odot\le r\le r_{\mathrm{out}}$.
Скорость направлена вдоль поля; все поля зависят только от $r,t$.
Теплопроводность, радиационные потери и дополнительный нагрев отсутствуют (пока что).
Газ однотемпературный, $\gamma=5/3$ и средняя молекулярная масса
$\mu$ постоянны. Единственный источник тепла --- диссипация волн.

Искомые физические поля: плотность $\rho$, скорость $u$,
газовое давление $p$ и плотности энергий встречных альфвеновских
волн $w_\pm\ge0$. Выбираем поле направленным в сторону роста $r$:
\[
 B(r)=\frac{A(R_\odot)B_\ph}{A(r)},\qquad
 \VA=\frac{B}{\sqrt{\mu_0\rho}},\qquad
 \PA=\frac{w_++w_-}{2},\qquad
 p=\frac{\rho k_{\mathrm B}T}{\mu m_u}.
\]
Здесь $\mu_0$ --- магнитная постоянная, $m_u$ --- атомная единица массы.
Регулярное магнитное поле задаёт $\VA$ и геометрию;
его продольная сила Лоренца равна нулю при $\bm u\parallel\bm B$.

\begin{equation}
\begin{aligned}
 \frac{\partial\rho}{\partial t}+\frac1A\frac{\partial(A\rho u)}{\partial r}&=0,\\[3pt]
 \rho\left(\frac{\partial u}{\partial t}+u\frac{\partial u}{\partial r}\right)
 &=-\frac{\partial(p+\PA)}{\partial r}-\rho\frac{GM_\odot}{r^2},\\[3pt]
 \frac{\partial p}{\partial t}+u\frac{\partial p}{\partial r}+
 \frac{\gamma p}{A}\frac{\partial(Au)}{\partial r}&=(\gamma-1)Q_w,\\[3pt]
 \frac{\partial w_\pm}{\partial t}+\frac1A\frac{\partial}{\partial r}\!\left[A(u\pm\VA)w_\pm\right]
 +\frac{w_\pm}{2A}\frac{\partial(Au)}{\partial r}
 &=\mp\Rref\sqrt{w_+w_-}-\Gamma_\pm w_\pm.
\end{aligned}
\label{eq:primitive}
\end{equation}
Знак $+$ соответствует волне вдоль поля относительно плазмы,
знак $-$ --- против поля. В системе отсчёта Солнца скорости переноса
равны $u\pm\VA$.

\textbf{Турбулентное замыкание.}
\[
 L_\perp(r)=L_{\perp,\ph}\sqrt{\frac{A(r)}{A(R_\odot)}},\qquad
 \Gamma_\pm=\frac{2}{L_\perp}\sqrt{\frac{w_\mp}{\rho}},\qquad
 Q_w=\Gamma_+w_++\Gamma_-w_-.
\]
Используем то же отражение, что в исходном тексте:
\[
 \Rref=\mathcal C(w_+,w_-)
 \min\!\left\{\left|\frac{\partial\VA}{\partial r}\right|,
                 \max(\Gamma_+,\Gamma_-)\right\},
\]
\[
 \mathcal C=
 \begin{cases}
  1-2\sqrt{w_-/w_+},&w_+\ge4w_-,\quad w_+>0,\\[3pt]
  2\sqrt{w_+/w_-}-1,&w_-\ge4w_+,\quad w_->0,\\[3pt]
  0,&\text{в остальных случаях}.
 \end{cases}
\]
В частности, $\mathcal C(0,0)=0$.
Отражение перераспределяет энергию между волнами,
а диссипация передаёт её газу.

\section{Переход к консервативным переменным}

Плотность кинетической и внутренней энергии газа равна
\[
 H=\frac{\rho u^2}{2}+\frac{p}{\gamma-1}.
\]
Для гладких решений система \eqref{eq:primitive} даёт
\[
 \frac{\partial H}{\partial t}+\frac1A\frac{\partial}{\partial r}\!\left[Au(H+p)\right]
 =-\rho u\frac{GM_\odot}{r^2}-u\frac{\partial\PA}{\partial r}+Q_w.
\]
Обозначим $W=w_++w_-$ и сложим два волновых уравнения:
\[
 \frac{\partial W}{\partial t}+\frac1A\frac{\partial}{\partial r}
 \!\left[A\bigl(uW+\VA(w_+-w_-)\bigr)\right]
 =-\frac{\PA}{A}\frac{\partial(Au)}{\partial r}-Q_w.
\]
При сложении этих балансов $Q_w$ сокращается, как и отражение,
а работа давления образует дивергенцию:
\[
 u\frac{\partial\PA}{\partial r}+\frac{\PA}{A}\frac{\partial(Au)}{\partial r}
 =\frac1A\frac{\partial(Au\PA)}{\partial r}.
\]
Поэтому для полной энергии \emph{газа и волн}
\[
 E=H+w_++w_-=
 \frac{\rho u^2}{2}+\frac{p}{\gamma-1}+w_++w_-
\]
получаем консервативный баланс
\begin{equation}
 \frac{\partial(AE)}{\partial t}+\frac{\partial}{\partial r}
 \!\left[A\bigl(u(E+p+\PA)+\VA(w_+-w_-)\bigr)\right]
 =-A\rho u\frac{GM_\odot}{r^2}.
\label{eq:totalenergy}
\end{equation}
Нагрев не исчез: при уменьшении $w_++w_-$ и фиксированной полной
энергии возрастает внутренняя энергия газа. Добавлять $Q_w$
в \eqref{eq:totalenergy} повторно нельзя.

\section{Консервативно-балансная запись}

Используя \eqref{eq:primitive} и \eqref{eq:totalenergy}, получаем
\begin{equation}
 \frac{\partial(A\U)}{\partial t}
 +\frac{\partial(A\F)}{\partial r}
 +\bm K\frac{\partial(Au)}{\partial r}=A\Src,
 \label{eq:vector}
\end{equation}
где
\[
 \U=\begin{pmatrix}
 \rho\\ \rho u\\ E\\ w_+\\ w_-
 \end{pmatrix},\qquad
 \F=\begin{pmatrix}
 \rho u\\
 \rho u^2+p+\PA\\
 u(E+p+\PA)+\VA(w_+-w_-)\\
 (u+\VA)w_+\\
 (u-\VA)w_-
 \end{pmatrix},\qquad
 \bm K=\begin{pmatrix}
 0\\0\\0\\w_+/2\\w_-/2
 \end{pmatrix},
\]
\[
 \Src=\begin{pmatrix}
 0\\[2pt]
 (p+\PA)A'/A-\rho GM_\odot/r^2\\[2pt]
 -\rho uGM_\odot/r^2\\[2pt]
 -\Rref\sqrt{w_+w_-}-\Gamma_+w_+\\[2pt]
 \Rref\sqrt{w_+w_-}-\Gamma_-w_-
 \end{pmatrix},\qquad A'=\frac{\dd A}{\dd r}.
\]
Давление восстанавливается по полной энергии:
\[
 \PA=\frac{w_++w_-}{2},\qquad
 p=(\gamma-1)\left(E-\frac{\rho u^2}{2}-w_+-w_-\right).
\]
Допустимое состояние требует $\rho>0$, $p>0$, $w_\pm\ge0$.

Кроме члена сжатия, пространственную производную содержит отражение:
\[
 \frac{\partial\VA}{\partial r}=-\VA
 \left(\frac{A'}A+\frac1{2\rho}\frac{\partial\rho}{\partial r}\right).
\]
Поэтому в \eqref{eq:vector} точно следует понимать
$\Src=\Src(\U,r,\partial\rho/\partial r)$.
На градиентной ветви отражение также является дифференциальной связью.


\section{Обезразмеривание}

Сохраняем гравитационную нормировку исходного текста:
\[
 R_*=R_\odot,\qquad
 U_*=\sqrt{\frac{GM_\odot}{R_*}},\qquad
 \rho_*=\rho_\ph,\qquad T_*=T_\ph,\qquad A_*=A(R_*).
\]
Масштабы давления, энергий и времени выбираем согласованно:
\begin{center}
\small
\setlength{\tabcolsep}{8pt}
\begin{tabular}{@{}ll@{\qquad}ll@{}}
\toprule
Величина & Масштаб & Величина & Масштаб\\
\midrule
Длина & $R_*$ & Время & $R_*/U_*$\\
Скорость & $U_*$ & Плотность & $\rho_*$\\
Температура & $T_*$ & Площадь & $A_*$\\
Поле & $B_*=B_\ph$ & Давление, энергии & $p_*=\rho_*U_*^2$\\
Частоты & $U_*/R_*$ & Нагрев & $p_*U_*/R_*$\\
\bottomrule
\end{tabular}
\end{center}

Вводим $\hat r=r/R_*$, $\hat t=t/t_*$,
$\hat\rho=\rho/\rho_*$, $\hat u=u/U_*$ и аналогично делим
остальные величины на масштабы из таблицы.
Далее шапочки опускаем: все поля и координаты ниже безразмерны.
Тогда $A(1)=1$, $B=1/A$, а коэффициент гравитации равен единице.

\subsection*{Безразмерные коэффициенты}

При фиксированных $\gamma,\mu$ используем три коэффициента:
\begin{equation}
 \theta=\frac{k_{\mathrm B}T_*}{\mu m_uU_*^2},\qquad
 b=\frac{B_*}{U_*\sqrt{\mu_0\rho_*}},\qquad
 d=\frac{2R_*}{L_{\perp,\ph}}.
 \label{eq:parameters}
\end{equation}
Они характеризуют соответственно тепловую энергию относительно
гравитационной, альфвеновскую скорость относительно $U_*$
и отношение глобального масштаба к корреляционной длине.
В частности, $p_\ph=\theta p_*$, а $V_{A,\ph}=bU_*$.

При желании первые два коэффициента можно выразить через
масштабные числа Маха: $\theta=1/(\gamma M_*^2)$ и $b=1/M_{A,*}$,
где $M_*=U_*/\sqrt{\gamma k_{\mathrm B}T_*/(\mu m_u)}$,
$M_{A,*}=U_*/V_{A,\ph}$. Далее эти обозначения не используются.

\section{Итоговая безразмерная система}

Используем исходные волновые энергии $w_\pm$ и полную энергию $E$:
масса, импульс и энергия записаны в консервативном виде,
а в волновых балансах явно сохранён член сжатия:
\begin{equation}
\begin{aligned}
 \frac{\partial(A\rho)}{\partial t}+\frac{\partial(A\rho u)}{\partial r}&=0,\\[4pt]
 \frac{\partial(A\rho u)}{\partial t}+\frac{\partial}{\partial r}\!\left[A(\rho u^2+p+\PA)\right]
 &=(p+\PA)A'-\frac{A\rho}{r^2},\\[4pt]
 \frac{\partial(AE)}{\partial t}+\frac{\partial}{\partial r}
 \!\left[A\bigl(u(E+p+\PA)+\VA(w_+-w_-)\bigr)\right]
 &=-\frac{A\rho u}{r^2},\\[4pt]
 \frac{\partial(Aw_+)}{\partial t}
 +\frac{\partial}{\partial r}\!\left[A(u+\VA)w_+\right]
 +\frac{w_+}{2}\frac{\partial(Au)}{\partial r}
 &=-A\left(\Rref\sqrt{w_+w_-}+\Gamma_+w_+\right),\\[4pt]
 \frac{\partial(Aw_-)}{\partial t}
 +\frac{\partial}{\partial r}\!\left[A(u-\VA)w_-\right]
 +\frac{w_-}{2}\frac{\partial(Au)}{\partial r}
 &=A\left(\Rref\sqrt{w_+w_-}-\Gamma_-w_-\right).
\end{aligned}
\label{eq:dimensionless}
\end{equation}
Система \eqref{eq:dimensionless} замыкается соотношениями
\[
\begin{gathered}
 \PA=\frac{w_++w_-}{2},\qquad
 p=(\gamma-1)\left(E-\frac{\rho u^2}{2}-w_+-w_-\right),
 \qquad T=\frac{p}{\theta\rho},\\[4pt]
 \VA=\frac{b}{A\sqrt\rho},\qquad
 \Gamma_\pm=d\sqrt{\frac{w_\mp}{A\rho}},\qquad
 Q_w=\Gamma_+w_++\Gamma_-w_-,\\[4pt]
 \Rref=\mathcal C(w_+,w_-)
 \min\!\left\{
 b\left|\frac{\partial}{\partial r}\frac1{A\sqrt\rho}\right|,
 d\sqrt{\frac{\max(w_+,w_-)}{A\rho}}
 \right\}.
\end{gathered}
\]
Функция $\mathcal C$ остаётся прежней, поскольку зависит
только от отношения энергий волн. Здесь $Q_w$ --- диагностическая
скорость нагрева; в уравнение полной энергии она не добавляется.

\subsection*{Куда вошли коэффициенты}

В системе \eqref{eq:dimensionless} коэффициенты входят через замыкания:

\textbf{$\theta$: уравнение состояния и тепловые данные.}
Связь $p=\theta\rho T$ задаёт пересчёт температуры, а на нижней
границе при $\rho=T=1$ получаем $p=\theta$.
При расчёте через $E$ давление восстанавливается из энергии,
поэтому $\theta$ не стоит явно в потоках: при отсутствии зависящих
от температуры источников оно входит в начальные и граничные
тепловые данные и в формулу $T=p/(\theta\rho)$.
Если вместо $p$ подставить $\theta\rho T$, его вклад в поток
импульса и энергии становится явным.

\textbf{$b$: перенос и отражение волн.}
Через $V_A=b/(A\sqrt\rho)$ он входит в скорости переноса $u\pm V_A$,
в поток полной энергии $V_A(w_+-w_-)$ и в градиентную ветвь отражения.
Газ испытывает его влияние через изменение волн и волнового давления.

\textbf{$d$: диссипация и нагрев.}
Он умножает частоты $\Gamma_\pm$, поэтому определяет скорость
убывания волновой энергии и её перехода в тепло.
Он также входит в ограничение отражения через $\max(\Gamma_+,\Gamma_-)$.
В уравнении полной энергии отдельного источника $Q_w$ нет:
обмен энергией газа и волн является внутренним и сокращается.
Коэффициент гравитации равен единице благодаря выбору $U_*^2=GM_\odot/R_*$.

\textbf{Данные задачи.}
Помимо коэффициентов \eqref{eq:parameters}, задаются геометрия $A(r)$,
размер области $r_{\mathrm{out}}/R_*$, начальные профили всех пяти полей
и граничные данные для входящих характеристик связанной системы.
Например, нормировка нижней границы даёт
\[
 \rho(1,t)=1,\qquad T(1,t)=1,\qquad p(1,t)=\theta,\qquad
 w_+(1,t)=\omega(t).
\]
\[
 \omega=\frac{w_{+,\ph}}{\rho_*U_*^2}
        =\left(\frac{\delta v_{+,\ph}}{U_*}\right)^2,
 \qquad w_{+,\ph}=\rho_*\delta v_{+,\ph}^{\,2}.
\]
Здесь $\delta v_{+,\ph}$ --- среднеквадратичная скорость
отдельной выходящей популяции волн. Набор граничных условий
должен соответствовать числу входящих характеристик;
значения $u$ и $w_-$ нельзя автоматически фиксировать вместе
с перечисленными данными. Перечисленные равенства не заменяют
полного характеристического анализа.
Массовый расход $A\rho u$ постоянен по $r$ только в стационарном режиме.

\section{Характерные значения параметров}

Для опорного примера: $\gamma=5/3$, $\mu=0.6$,
$R_*=695700$ км, $U_*\simeq436.76$ км/с.
Постоянная $\mu$ не описывает изменение ионизации по высоте.
Ниже приведены опорные значения и иллюстративные диапазоны
для проверки чувствительности модели, а не наблюдательные интервалы.
\begin{center}
\small
\setlength{\tabcolsep}{9pt}
\begin{tabular}{@{}lrl@{}}
\toprule
Величина & Опорное значение & Диапазон перебора\\
\midrule
$\rho_*$, кг/м$^3$ & $0.0003$ & $0.0002$ -- $0.0004$\\
$T_*$, К & $6400$ & $5000$ -- $7000$\\
$B_*$, Тл & $0.0001$ & $0.0001$ -- $0.001$\\
$L_{\perp,\ph}$, км & $100$ & $50$ -- $200$\\
\midrule
$\theta$ & $0.000465$ & $0.000363$ -- $0.000509$\\
$b$ & $0.0000118$ & $0.0000102$ -- $0.000144$\\
$d$ & $13914$ & $6957$ -- $27828$\\
\bottomrule
\end{tabular}
\end{center}
Диапазоны коэффициентов \eqref{eq:parameters} рассчитаны по крайним
значениям исходных данных при независимом их варьировании.
При опорных данных $t_*\simeq1593$ с,
$p_*\simeq57.23$ МПа, $Q_*\simeq35928$ Вт/м$^3$,
$V_{A,\ph}\simeq5.15$ м/с.

Помимо $\theta,b,d$, задаются геометрия, начальные профили,
граничная волновая энергия $\omega$ и размер области.
Для внешней границы на 1 а.е. $r_{\mathrm{out}}/R_*=215.03$.
Для малых альфвеновских возмущений требуется
\[
 \frac{\delta v_{+,\ph}}{V_{A,\ph}}=\frac{\sqrt\omega}{b}\ll1,
 \qquad\text{то есть}\qquad \omega\ll b^2.
\]
При $B_*=0.0001$ Тл и радиальной трубке $A=r^2$ поле на 1 а.е.
равно примерно $2.16$ нТл. Изменение поля основания требует
согласовать геометрию и внешнее поле; таблица диапазонов
не заменяет этой проверки.

\newpage
\section{Характеристики и граничные условия}

Рассматриваем безразмерную модель без теплопроводности и радиационных
потерь на $1\le r\le r_{\rm out}$, с заданной стационарной площадью $A(r)$.
Анализ локальный: коэффициенты замораживаются в рассматриваемом
гладком состоянии. Предполагаем $\rho>0$, $p>0$, $w_\pm>0$. Пять неизвестных дают характеристический полином пятой степени;
существование пяти вещественных характеристик ещё нужно проверить.

\subsection{Главная часть системы}

Для анализа удобно использовать исходные физические переменные
\[
 \bm q=(\rho,u,p,w_+,w_-)^{\mathsf T},\qquad
 c_s^2=\frac{\gamma p}{\rho},\qquad
 \VA=\frac{b}{A\sqrt\rho},\qquad a=\frac{A'}A.
\]
Переход от $(\rho,\rho u,E,w_+,w_-)$ к $\bm q$ обратим, поэтому характеристические скорости не меняются.
Все пространственные производные неизвестных необходимо включить
в главную часть, в том числе производную в отражении.

Обозначим $G=\max(\Gamma_+,\Gamma_-)$ и
$\sigma=\operatorname{sign}(\partial\VA/\partial r)$.
Замыкание отражения имеет две гладкие ветви:
\[
 \Rref=\begin{cases}
 \mathcal C\sigma\,\partial\VA/\partial r,
     &0<|\partial\VA/\partial r|<G,\\
 \mathcal C G,&|\partial\VA/\partial r|>G.
 \end{cases}
\]
При $\mathcal C=0$ отражение отсутствует. Для компактности введём один коэффициент
\begin{equation}
 h=\begin{cases}
 \mathcal C\sigma\sqrt{w_+w_-},
     &0<|\partial\VA/\partial r|<G,\\
 0,&|\partial\VA/\partial r|>G\ \text{или}\ \mathcal C=0.
 \end{cases}
 \label{ch:h}
\end{equation}
Коэффициент \eqref{ch:h} не является новой искомой переменной. Это, фактически, <<переключатель>>, он включает или выключает дополнительный вклад производной, связанный с отражением. Если отражения нет, то $h = 0$, поэтому слева новой производной не добавится. Если есть, то производная будет.

Из сохранения магнитного потока следует (хотим от производной по альфвеновской скорости перейти к производной плотности):
\[
 \frac{\partial\VA}{\partial r}
 =-a\VA-\frac{\VA}{2\rho}\frac{\partial\rho}{\partial r}.
\]
Например, уравнение для $w_+$ до переноса градиентного отражения влево (так мы полностью все производные переносим влево):
\[
 \frac{\partial w_+}{\partial t}
 +(u+\VA)\frac{\partial w_+}{\partial r}
 +\frac32w_+\frac{\partial u}{\partial r}
 -\frac{\VA w_+}{2\rho}\frac{\partial\rho}{\partial r}
 =-\frac32au w_+-\Rref\sqrt{w_+w_-}-\Gamma_+w_+.
\]
Поэтому система на фиксированной ветви записывается как
\begin{equation}
 \frac{\partial\bm q}{\partial t}
 +\mathsf J\frac{\partial\bm q}{\partial r}=\bm s,
 \qquad
 \mathsf J=\begin{pmatrix}
 u&\rho&0&0&0\\
 0&u&1/\rho&1/(2\rho)&1/(2\rho)\\
 0&\gamma p&u&0&0\\
 -\VA(w_++h)/(2\rho)&3w_+/2&0&u+\VA&0\\
 \VA(w_-+h)/(2\rho)&3w_-/2&0&0&u-\VA
 \end{pmatrix}.
 \label{ch:matrix}
\end{equation}
Матрица \eqref{ch:matrix} -- матрица коэффициентов при пространственных производных. Источник:
\[
 \bm s=\begin{pmatrix}
 -a\rho u\\ -1/r^2\\
 -\gamma aup+(\gamma-1)Q_w\\
 -3au w_+/2-J_0-\Gamma_+w_+\\
 -3au w_-/2+J_0-\Gamma_-w_-
 \end{pmatrix},\qquad
J_0=
\begin{cases}
-aV_Ah,
& 0<\left|\dfrac{\partial V_A}{\partial r}\right|<G,\\[8pt]
\mathcal C G\sqrt{w_+w_-},
& \left|\dfrac{\partial V_A}{\partial r}\right|>G,\\[8pt]
0,
& \mathcal C=0.
\end{cases}
\]

Итак, смысл системы \eqref{ch:matrix}: если есть отражение, то производная уже стоит слева, а справа их нет. Если же отражения нет, то ни слева, ни справа не будет дополнительной производной.

\subsection{Характеристический полином и пять скоростей}

Характеристические скорости определяются уравнением
$\det(\lambda\mathsf I-\mathsf J)=0$.
Положим $z=\lambda-u$, $W=w_++w_-$, $\Delta W=w_+-w_-$.
Для возмущения $(\delta\rho,\delta u,\delta p,\delta w_+,\delta w_-)$
при $z\ne0,\pm\VA$ первое и третье уравнения дают
\[
 \delta\rho=\frac{\rho}{z}\delta u,\qquad
 \delta p=\frac{\gamma p}{z}\delta u,
 \qquad
 \delta w_\pm=
 \frac{3z w_\pm\mp\VA(w_\pm+h)}{2z(z\mp\VA)}\delta u.
\]
Подставляя их в уравнение импульса и раскрывая определитель,
получаем:
\begin{equation}
 \det(\lambda\mathsf I-\mathsf J)=zP(z),\qquad
 P(z)=(z^2-c_s^2)(z^2-\VA^2)
 -\frac{3Wz^2+2\VA\Delta Wz-\VA^2(W+2h)}{4\rho}.
 \label{ch:poly}
\end{equation}
Таким образом, одна скорость точно равна $u$, остальные четыре равны
$u+z_j$, где $z_j$ --- корни квартета \eqref{ch:poly}.
Для вычисления достаточно найти корни полинома четвёртой степени
или собственные значения матрицы; формула в радикалах здесь не нужна.

При исчезающе малых энергиях волн полином стремится к
$(z^2-c_s^2)(z^2-\VA^2)$, и скорости стремятся к
\[
 u,\qquad u\pm c_s,\qquad u\pm\VA.
\]
При конечных $w_\pm$ это только приближение.
Волновое давление меняет скорости, а на градиентной ветви их
дополнительно меняет $h$. Геометрия, гравитация и локальная диссипация
не дают отдельных характеристик. Параметр $d$ влияет на выбор ветви
отражения, а $b$ входит в матрицу через $\VA$.

Точное решение для этого уравнения четвертой степени можно найти, однако формулы там получаются громоздкие. Решение дано в Приложении \ref{app:characteristic_roots}. Однако и оно получено в предположении, что существуют пять различных вещественных корней.

\textbf{Подумать: оценить, может отношение W к cs cs rho и к rho Va Va? Тогда станет понятно, можно ли пренебрегать}

\subsection{Проверка гиперболичности}

\textbf{Ограниченная ветвь и отсутствие отражения.}
Здесь $h=0$. Из \eqref{ch:poly} следует
\[
 P(0)=\VA^2\left(c_s^2+\frac W{4\rho}\right)>0,
 \qquad P(\VA)=-\frac{\VA^2w_+}{\rho}<0,
 \qquad P(-\VA)=-\frac{\VA^2w_-}{\rho}<0.
\]
Поскольку $P(z)\to+\infty$ при $z\to\pm\infty$, в каждом
из четырёх интервалов имеется корень. Степень полинома равна четырём,
поэтому эти корни единственны, вещественны и просты:
\begin{equation}
 z_1<-\VA<z_2<0<z_3<\VA<z_4.
 \label{ch:order}
\end{equation}
Следовательно, система строго гиперболична, и её упорядоченные скорости
\begin{equation}
 \lambda_1=u+z_1,\quad\lambda_2=u+z_2,\quad
 \lambda_3=u,\quad\lambda_4=u+z_3,\quad\lambda_5=u+z_4.
 \label{ch:speeds}
\end{equation}

\textbf{Градиентная ветвь.}
Здесь
\[
 P(0)=\VA^2\left(c_s^2+\frac{W+2h}{4\rho}\right)>0,
 \qquad
 P(\pm\VA)=-\frac{\VA^2}{\rho}\left(w_\pm-\frac h2\right).
\]
Первое неравенство следует из $|\mathcal C|\le1$,
$|h|\le\sqrt{w_+w_-}$ и $W\ge2\sqrt{w_+w_-}$.
Поэтому достаточное условие строгой гиперболичности и порядка корней
\eqref{ch:order} имеет вид
\begin{equation}
 h<2\min(w_+,w_-).
 \label{ch:sufficient}
\end{equation}

Действительно, в таком случае будет четыре смены знака, а значит, четыре действительных корня: 

\begin{center}
\begin{tikzpicture}[>=Stealth, x=1.5cm]
  % ось z
  \draw[->] (-4.4,0) -- (4.6,0) node[right] {$z$};

  % отмеченные точки и подписи под осью
  \foreach \x/\lab in {-4/-\infty, -2/-V_A, 0/0, 2/V_A, 4/+\infty}
    \draw (\x,0.08) -- (\x,-0.08) node[below=2pt] {$\lab$};

  % знак P в отмеченных точках (над осью)
  \foreach \x/\s in {-4/+, -2/-, 0/+, 2/-, 4/+}
    \node[circle, draw, inner sep=1.5pt, above=5pt] at (\x,0) {$\s$};

  % корни в каждом из четырёх интервалов
  \foreach \x/\n in {-3/1, -1/2, 1/3, 3/4}
    \node[red] at (\x,0) {$\times$};
  \foreach \x/\n in {-3/1, -1/2, 1/3, 3/4}
    \node[red, below=14pt] at (\x,0) {$z_\n$};

  % подпись: смена знака на каждом интервале
  \foreach \x in {-3,-1,1,3}
    \draw[<->, gray] (\x-0.7,0.9) -- (\x+0.7,0.9)
      node[midway, above] {\footnotesize смена знака};
\end{tikzpicture}
\end{center}

В частности, это выполняется при $h\le0$.
Условие \eqref{ch:sufficient} достаточное, но не необходимое:
если оно нарушено, нужно непосредственно проверить корни.

\emph{Для данного замыкания гиперболичность не гарантирована глобально.}
Контрпример: $\rho=1$, $c_s=\VA=1$, $w_+=0.1$, $w_-=0.001$,
$\mathcal C=0.8$, $\sigma=1$, откуда $h=0.008$.
Тогда
\[
 P(z)=z^4-2.07575z^2-0.0495z+1.02925,
\]
\[
 z\simeq0.858576,\quad1.167488,\quad-1.013032\pm0.023976\,\mathrm i.
\]
Это допустимое состояние градиентной ветви: например,
при $r=A=b=1$, $A'/A=2$, $d=1$, $\partial\rho/\partial r=-4.2$
имеем $\partial\VA/\partial r=0.1<G=\sqrt{0.1}$.
Две скорости комплексные, и обычная постановка через пять вещественных
входящих/выходящих характеристик здесь неприменима.
Замораживание отражения как источника при расщеплении даёт матрицу
с $h=0$ для подшага, но не доказывает гиперболичность полной модели.

\subsection{Сколько характеристик входит через каждую границу}

Далее предполагаем гиперболическое состояние с двумя отрицательными
и двумя положительными $z_j$, упорядоченными как $z_1<z_2<0<z_3<z_4$,
и истечение $u>0$. На ограниченной ветви это доказано выше;
на градиентной проверяется по корням.
Характеристика удовлетворяет $\mathrm dr/\mathrm dt=\lambda_j$.
Для скоростей \eqref{ch:speeds} слева $\lambda>0$ означает вход {из фотосферы в область},
а $\lambda<0$ --- выход из области к фотосфере.
Справа направления относительно области меняются местами.

\begin{center}
\small
\begin{tabular}{@{}lcc@{}}
\toprule
Режим в данной граничной точке & Слева: вход / выход & Справа: вход / выход\\
\midrule
$0<u<-z_2$ & $3/2$ & $2/3$\\
$-z_2<u<-z_1$ & $4/1$ & $1/4$\\
$u>-z_1$ & $5/0$ & $0/5$\\
\bottomrule
\end{tabular}
\end{center}
Каждый столбец применяется к состоянию на своей границе.

\textbf{Фотосфера.}
На нижней границе предполагается дозвуковое и
доальфвеновское истечение:
\[
0<u_\ph<\min(c_{s,\ph},V_{A,\ph}).
\]
Оценим справедливость этого предположения.
При характерных параметрах фотосферы
$\rho_\ph=3\cdot10^{-4}$ кг/м$^3$,
$T_\ph=6400$ К, $B_\ph=10^{-4}$ Тл,
$\gamma=5/3$, $\mu=1.3$ получаем
\[
c_{s,\ph}=
\sqrt{\frac{\gamma k_{\mathrm B}T_\ph}{\mu m_u}}
\simeq 8.3\ \text{км/с},
\qquad
V_{A,\ph}=
\frac{B_\ph}{\sqrt{\mu_0\rho_\ph}}
\simeq 5.15\ \text{м/с}.
\]

Для оценки скорости истечения воспользуемся
типичными параметрами солнечного ветра на 1 а.е.:
$n_{\mathrm{AU}}\simeq5$ см$^{-3}$,
$u_{\mathrm{AU}}\simeq400$ км/с.
Из сохранения массового расхода в стационарной
радиальной трубке $A\propto r^2$ следует
\[
u_\ph=
\frac{m_p n_{\mathrm{AU}}u_{\mathrm{AU}}}{\rho_\ph}
\left(\frac{r_{\mathrm{AU}}}{R_\odot}\right)^2
\simeq 5.2\cdot10^{-7}\ \text{м/с}.
\]
Таким образом,
\[
\boxed{u_\ph\ll V_{A,\ph}\ll c_{s,\ph}.}
\]
Следовательно, в приближении слабых волн
на фотосфере три характеристики входят
в расчётную область, а две выходят из неё.

\textbf{Дальняя граница.}
Условия извне не нужны, так как на правой границе течение сверхзвуковое и сверхальфвеновское. Все характеристики выходят из области течения.

\subsection{Выходящие характеристики}

На левой границе рассматриваем режим, когда три характеристики
входят в расчётную область, а две выходят из неё к фотосфере.
Две выходящие характеристики приносят информацию
из внутреннего решения и позволяют определить
скорость $u_{\mathrm{ph}}$ и энергию встречных волн
$w_{-,\mathrm{ph}}$.

Обозначим столбец искомых величин через
\[
 \bm q=(\rho,u,p,w_+,w_-)^{\mathsf T}.
\]
Система имеет вид
\[
\frac{\partial\bm q}{\partial t}
+\mathsf J\frac{\partial\bm q}{\partial r}
=\bm s.
\]
Умножим её слева на строку коэффициентов
$\bm\ell^{\mathsf T}$:
\[
\bm\ell^{\mathsf T}\frac{\partial\bm q}{\partial t}
+\bm\ell^{\mathsf T}\mathsf J
\frac{\partial\bm q}{\partial r}
=\bm\ell^{\mathsf T}\bm s.
\]
Хотим, чтобы левая часть представляла собой производную
вдоль кривой $dr/dt=\lambda$, то есть
\[
\bm\ell^{\mathsf T}
\left(
\frac{\partial\bm q}{\partial t}
+\lambda\frac{\partial\bm q}{\partial r}
\right).
\]
Для этого необходимо потребовать
\[
\boxed{
\bm\ell^{\mathsf T}\mathsf J
=\lambda\bm\ell^{\mathsf T}.
}
\]
Это определение левого собственного вектора матрицы
$\mathsf J$, а $\lambda$ --- соответствующее собственное значение.
Таким образом, вдоль характеристики $dr/dt=\lambda$
получаем условие
\[
\boxed{
\bm\ell^{\mathsf T}\frac{d\bm q}{dt}
=\bm\ell^{\mathsf T}\bm s,
\qquad
\frac{dr}{dt}=\lambda.
}
\]

\subsubsection{Соотношения вдоль выходящих траекторий.}

Для первого приближения пренебрежём влиянием
волнового давления на движение газа, нагревом
и геометрическим членом в уравнении давления.
Гравитацию сохраним.
Тогда
\[
 \partial_tu+u\partial_ru+\frac1\rho\partial_rp
 =-\frac1{r^2},
 \qquad
 \partial_tp+u\partial_rp+\rho c_s^2\partial_ru=0.
\]
Вычтем из второго уравнения первое,
умноженное на $\rho c_s$:
\[
 \bigl[\partial_t+(u-c_s)\partial_r\bigr]p
 -\rho c_s
 \bigl[\partial_t+(u-c_s)\partial_r\bigr]u
 =\frac{\rho c_s}{r^2}.
\]
Таким образом, вдоль выходящей акустической траектории
получаем
\begin{equation}
 \boxed{
 \frac{dp}{dt}-\rho c_s\frac{du}{dt}
 \simeq\frac{\rho c_s}{r^2},
 \qquad
 \frac{dr}{dt}=u-c_s.
 }
 \label{bc:acoustic-path}
\end{equation}
Здесь $d/dt$ означает производную именно вдоль
указанной траектории.
Для вывода не требовалось считать $\rho c_s$
постоянным; это дополнительное приближение
будет использовано при интегрировании на одном шаге.

Для встречных волн примем более грубое приближение:
\begin{equation}
 \boxed{
 \frac{dw_-}{dt}\simeq0,
 \qquad
 \frac{dr}{dt}=u-\VA.
 }
 \label{bc:wave-path}
\end{equation}
Оно допустимо, если за время $\Delta t$ отражение,
диссипация, сжатие и неоднородность среды
мало меняют $w_-$.
Слабое влияние волн на газ само по себе
этого не гарантирует.

\subsection{Граничные условия на фотосфере}

На новом временном слое задаём
\[
\rho_{\mathrm{ph}}^{n+1}=1,\qquad
T_{\mathrm{ph}}^{n+1}=1,\qquad
p_{\mathrm{ph}}^{n+1}=\theta,\qquad
w_{+,\mathrm{ph}}^{n+1}=\omega(t^{n+1}).
\]
Скорость $u_{\mathrm{ph}}^{n+1}$ и энергию встречных волн
$w_{-,\mathrm{ph}}^{n+1}$ определяем по информации
из расчётной области.
Ниже используем упрощённые акустическое и волновое
соотношения совместности.

\subsubsection{Основания характеристик}

На старом слое $t^n$ решение представлено постоянными
состояниями внутри ячеек.
Первая ячейка занимает интервал
\[
I_1=[r_{\mathrm{ph}},r_{3/2}],
\]
а $r_1$ обозначает её центр.
По известному состоянию этой ячейки вычисляем
\[
c_{s,1}^n=\sqrt{\frac{\gamma p_1^n}{\rho_1^n}},
\qquad
V_{A,1}^n=\frac{b}{A(r_1)\sqrt{\rho_1^n}}.
\]
На один шаг замораживаем характеристические скорости:
\[
\lambda_{\mathrm a}^n=u_1^n-c_{s,1}^n<0,
\qquad
\lambda_{\mathrm w}^n=u_1^n-V_{A,1}^n<0.
\]
Основания характеристик, приходящих в точку
$(r_{\mathrm{ph}},t^{n+1})$, находятся на старом слое:
\[
r_{\mathrm a}=r_{\mathrm{ph}}-\lambda_{\mathrm a}^n\Delta t,
\qquad
r_{\mathrm w}=r_{\mathrm{ph}}-\lambda_{\mathrm w}^n\Delta t.
\]
Выбираем шаг так, чтобы обе точки лежали в первой ячейке (в численном методе это условие автоматически будет выполнено -- условие Куранта-Фридрихса-Леви)ф:
\[
\max(-\lambda_{\mathrm a}^n,-\lambda_{\mathrm w}^n)\Delta t
\le r_{3/2}-r_{\mathrm{ph}}.
\]
Тогда вследствие кусочно-постоянного представления
\[
p_{\mathrm a}^n=p_1^n,\qquad
u_{\mathrm a}^n=u_1^n,\qquad
w_{-,\mathrm w}^n=w_{-,1}^n.
\]
Интерполяция не требуется.

% В преамбуле:
% \usepackage{tikz}
% \usetikzlibrary{arrows.meta}

\begin{figure}[htbp]
\centering
\begin{tikzpicture}[
    x=1.25cm,
    y=1.05cm,
    >=Stealth,
    font=\small
]
% Первая ячейка: от x=0 до x=6.
% Её центр: x=3.
% Основания характеристик: x=1.5 и x=4.5.

% Временные слои.
\draw[gray] (0,0) -- (6.4,0);
\draw[gray] (0,2.7) -- (6.4,2.7);
\node[left] at (-0.15,0) {$t^n$};
\node[left] at (-0.15,2.7) {$t^{n+1}$};

% Границы и центр первой ячейки.
\draw[thick] (0,-0.15) -- (0,3);
\draw[gray,dashed] (6,-0.15) -- (6,3);
\draw[gray,dotted] (3,0) -- (3,0.7);

\node[above] at (0,3.05) {фотосфера};
\node[above] at (6,3.05) {правая грань};

% Постоянное состояние на старом слое.
\draw[very thick,gray!65] (0,0) -- (6,0);
\fill (3,0) circle (1.8pt);
\node[above] at (3,0.7)
    {$\rho_1^n,\ u_1^n,\ p_1^n,\ w_{\pm,1}^n$};

% Характеристики.
\draw[blue!75!black,very thick,->]
    (1.5,0) -- (0,2.7);
\draw[red!75!black,very thick,->]
    (4.5,0) -- (0,2.7);

% Основания и новая граничная точка.
\fill[blue!75!black] (1.5,0) circle (2.4pt);
\fill[red!75!black] (4.5,0) circle (2.4pt);
\fill (0,2.7) circle (2.5pt);

\node[above right] at (0.1,2.7)
    {новое граничное состояние};

% Подписи координат.
\node[below] at (0,-0.12) {$r_{\mathrm{ph}}$};
\node[below,blue!75!black] at (1.5,-0.12)
    {$r_{\mathrm a}$};
\node[below] at (3,-0.12) {$r_1$};
\node[below,red!75!black] at (4.5,-0.12)
    {$r_{\mathrm w}$};
\node[below] at (6,-0.12) {$r_{3/2}$};

% Обозначения скоростей.
\node[blue!75!black,anchor=west] at (3.4,2.1)
    {$\lambda_{\mathrm a}^n=u_1^n-c_{s,1}^n$};
\node[red!75!black,anchor=west] at (3.4,1.6)
    {$\lambda_{\mathrm w}^n=u_1^n-V_{A,1}^n$};

% Размер первой ячейки.
\draw[<->] (0,-0.9) -- (6,-0.9)
    node[midway,below] {первая ячейка $I_1$};

\end{tikzpicture}
\caption{
Оба основания характеристик находятся в первой ячейке.
На старом слое состояние внутри неё постоянно,
поэтому значения в основаниях совпадают
с ячеечными значениями.
Взаимное расположение оснований показано условно.
}
\label{fig:characteristic_feet_constant}
\end{figure}

\subsubsection{Обновление граничного состояния}

Интегрируем соотношение~\eqref{bc:acoustic-path}
вдоль акустической характеристики.
Коэффициент $\rho c_s$ заменяем постоянным
значением $\rho_1^n c_{s,1}^n$ со старого слоя,
а в гравитационном источнике полагаем
$r\simeq r_{\mathrm{ph}}$:
\[
\theta-p_1^n
-\rho_1^n c_{s,1}^n
\left(u_{\mathrm{ph}}^{n+1}-u_1^n\right)
\simeq
\Delta t\,\frac{\rho_1^n c_{s,1}^n}{r_{\mathrm{ph}}^2}.
\]
В используемом упрощённом волновом соотношении
$w_-$ сохраняется вдоль характеристики.
Поэтому
\begin{equation}
\begin{aligned}
u_{\mathrm{ph}}^{n+1}
&\simeq
u_1^n+\frac{\theta-p_1^n}{\rho_1^n c_{s,1}^n}
-\frac{\Delta t}{r_{\mathrm{ph}}^2},\\[4pt]
w_{-,\mathrm{ph}}^{n+1}
&\simeq w_{-,1}^n.
\end{aligned}
\label{ch:leftsimple}
\end{equation}
Все величины справа известны на слое $t^n$.
В нашей нормировке $r_{\mathrm{ph}}=1$.

Полную энергию восстанавливаем после обновления:
\[
E_{\mathrm{ph}}^{n+1}
=
\frac{(u_{\mathrm{ph}}^{n+1})^2}{2}
+\frac{\theta}{\gamma-1}
+\omega(t^{n+1})
+w_{-,\mathrm{ph}}^{n+1}.
\]

\subsubsection{Возможность простого копирования}

Можно использовать более простую экстраполяцию:
\[
u_{\mathrm{ph}}^{n+1}=u_1^n,\qquad
w_{-,\mathrm{ph}}^{n+1}=w_{-,1}^n.
\]
Однако копирование скорости не учитывает её связь
с заданным давлением $\theta$ и гравитацией,
присутствующую в~\eqref{ch:leftsimple}.
Поэтому оно может искажать отражение акустических
возмущений и граничный массовый расход.
Для встречных волн копирование уже совпадает
с принятым здесь упрощённым характеристическим условием.

Схема~\eqref{ch:leftsimple} также является приближённой:
для полной связанной системы следует использовать
её точные выходящие характеристики и условия
совместности со всеми источниками.
Кусочно-постоянное представление само по себе
не обеспечивает точного сохранения
гидростатического равновесия.
\subsection{Граничные условия на выходе трубки}

Если на правой границе
\[
 \min_j\lambda_j(r_{\mathrm{out}},t)>0,
\]
то все пять характеристик выходят из области.
При обратном прослеживании каждая из них приходит
из внутренней точки предыдущего временного слоя.
Поэтому никаких дополнительных физических данных
от внешней среды здесь не требуется.

Для первого приближения используем копирование
последней расчётной ячейки:
\[
 \rho_{\mathrm{out}}=\rho_N,\qquad
 u_{\mathrm{out}}=u_N,\qquad
 p_{\mathrm{out}}=p_N,\qquad
 w_{\pm,\mathrm{out}}=w_{\pm,N}.
\]

Полную энергию вычисляем по тем же значениям:
\[
 E_{\mathrm{out}}
 =\frac{\rho_{\mathrm{out}}u_{\mathrm{out}}^2}{2}
  +\frac{p_{\mathrm{out}}}{\gamma-1}
  +w_{+,\mathrm{out}}+w_{-,\mathrm{out}}.
\]
Дополнительно фиксировать давление, скорость
или полагать $w_{-,\mathrm{out}}=0$ не следует.

\newpage

\section{Численный метод}


\subsection{Дивергентная часть и поток HLL или LF}

Безразмерная система \eqref{eq:dimensionless} записывается в виде
\begin{equation}
\boxed{
\frac{\partial(A\U)}{\partial t}
+\frac{\partial(A\F)}{\partial r}
=A\Src-\bm K\frac{\partial(Au)}{\partial r},
}
\end{equation}
где
\[
\U=
\begin{pmatrix}
\rho\\
\rho u\\
E\\
w_+\\
w_-
\end{pmatrix},
\qquad
\F=
\begin{pmatrix}
\rho u\\
\rho u^2+p+\PA\\
u(E+p+\PA)+V_A(w_+-w_-)\\
(u+V_A)w_+\\
(u-V_A)w_-
\end{pmatrix},
\]
\[
\Src=
\begin{pmatrix}
0\\[2pt]
(p+\PA)\dfrac{A'}{A}-\dfrac{\rho}{r^2}\\[6pt]
-\dfrac{\rho u}{r^2}\\[6pt]
-\Rref\sqrt{w_+w_-}-\Gamma_+w_+\\[4pt]
+\Rref\sqrt{w_+w_-}-\Gamma_-w_-
\end{pmatrix},
\qquad
\bm K=
\begin{pmatrix}
0\\0\\0\\w_+/2\\w_-/2
\end{pmatrix}.
\]
Здесь используются замыкания
\[
\PA=\frac{w_++w_-}{2},
\qquad
V_A=\frac{b}{A\sqrt{\rho}},
\qquad
p=(\gamma-1)
\left(E-\frac{\rho u^2}{2}-w_+-w_-\right).
\]


\subsubsection{Конечно-объёмная аппроксимация.}

Сначала рассмотрим только дивергентную часть системы,
формально положив $\Src=0$, $\bm K=0$:
\begin{equation}
\frac{\partial(A\U)}{\partial t}
+\frac{\partial(A\F)}{\partial r}=0.
\label{eq:homogeneous_system}
\end{equation}

Введём пространственную сетку с гранями
$r_{i-1/2}$, $r_{i+1/2}$ и временные слои
$t^n$, $t^{n+1}=t^n+\Delta t$.

Проинтегрируем \eqref{eq:homogeneous_system}
по пространственно-временной ячейке
$[r_{i-1/2},r_{i+1/2}]
\times[t^n,t^{n+1}]$:
\[
\int_{t^n}^{t^{n+1}}
\int_{r_{i-1/2}}^{r_{i+1/2}}
\left[
\frac{\partial(A\U)}{\partial t}
+\frac{\partial(A\F)}{\partial r}
\right]dr\,dt=0.
\]

Интегрируя первое слагаемое по времени,
а второе по пространству, получаем
\begin{align*}
&\int_{r_{i-1/2}}^{r_{i+1/2}}
A(r)\bigl[\U(r,t^{n+1})-\U(r,t^n)\bigr]\,dr
\\
&\quad+
\int_{t^n}^{t^{n+1}}
\left[
A_{i+1/2}\F(r_{i+1/2},t)
-A_{i-1/2}\F(r_{i-1/2},t)
\right]dt=0.
\end{align*}

Обозначим объём ячейки и среднее значение
консервативных переменных на временном слое:
\[
V_i=\int_{r_{i-1/2}}^{r_{i+1/2}}A(r)\,dr,
\qquad
\U_i^n=
\frac{1}{V_i}
\int_{r_{i-1/2}}^{r_{i+1/2}}
A(r)\U(r,t^n)\,dr.
\]

Кроме того, введём средний по времени
физический поток через грань:
\[
\overline{\F}_{i+1/2}^{\,n}
=\frac{1}{\Delta t}
\int_{t^n}^{t^{n+1}}
\F(r_{i+1/2},t)\,dt.
\]

Тогда предыдущее интегральное равенство
записывается точно:
\begin{equation}
\boxed{
\U_i^{n+1}
=\U_i^n
-\frac{\Delta t}{V_i}
\left(
A_{i+1/2}\overline{\F}_{i+1/2}^{\,n}
-A_{i-1/2}\overline{\F}_{i-1/2}^{\,n}
\right).
}
\label{eq:fv_exact}
\end{equation}

\subsubsection{Численные потоки.}

Формула \eqref{eq:fv_exact} пока точная,
но средние по времени потоки
$\overline{\F}_{i\pm1/2}^{\,n}$ неизвестны.
Их необходимо аппроксимировать.

Для схемы первого порядка используем
значения переменных на старом временном слое:
\[
\U_L=\U_i^n,\qquad
\U_R=\U_{i+1}^n.
\]
По этим состояниям вычисляем физические потоки
$\F_L$ и $\F_R$. Поскольку поток $\F$
зависит от альфвеновской скорости
$V_A=b/(A\sqrt{\rho})$, при его вычислении
используем
\[
V_{A,L}=\frac{b}{A_{i+1/2}\sqrt{\rho_L}},
\qquad
V_{A,R}=\frac{b}{A_{i+1/2}\sqrt{\rho_R}}.
\]
Площадь берётся на рассматриваемой грани,
а плотности --- из соответствующих состояний.

Вместо неизвестного среднего потока положим
\[
\overline{\F}_{i+1/2}^{\,n}
\approx\widehat{\F}_{i+1/2}^{\,n}.
\]

При наличии оценок крайних скоростей
$s_L<s_R$ поток HLL имеет вид
\begin{equation}
\widehat{\F}_{\mathrm{HLL}}=
\begin{cases}
\F_L, & s_L\ge0,\\[3pt]
\displaystyle
\frac{s_R\F_L-s_L\F_R+s_Ls_R(\U_R-\U_L)}
{s_R-s_L},
& s_L<0<s_R,\\[8pt]
\F_R, & s_R\le0.
\end{cases}
\label{eq:hll_flux}
\end{equation}

Более простой вариант --- локальный поток
Лакса--Фридрихса (Русанова):
\begin{equation}
\widehat{\F}_{\mathrm{LF}}
=\frac{\F_L+\F_R}{2}
-\frac{\alpha}{2}(\U_R-\U_L),
\label{eq:lf_flux}
\end{equation}
где $\alpha$ --- оценка максимального модуля
характеристической скорости.
Он соответствует выбору
$s_L=-\alpha$, $s_R=\alpha$ в HLL.

Подстановка численных потоков в
\eqref{eq:fv_exact} даёт окончательную формулу
схемы первого порядка для дивергентной части:
\begin{equation}
\boxed{
\U_i^{n+1}
=\U_i^n
-\frac{\Delta t}{V_i}
\left(
A_{i+1/2}\widehat{\F}_{i+1/2}^{\,n}
-A_{i-1/2}\widehat{\F}_{i-1/2}^{\,n}
\right).
}
\label{eq:fv_flux}
\end{equation}
Здесь $\widehat{\F}$ выбирается по формуле
\eqref{eq:hll_flux} или \eqref{eq:lf_flux}.

Полученная формула относится только
к дивергентной части системы.
Для исходных уравнений необходимо дополнительно
аппроксимировать источники $\Src$
и неконсервативный член
$\bm K\,\partial_r(Au)$.
Кроме того, требуется определить допустимые
оценки характеристических скоростей
с учётом возможной негиперболичности системы.



\subsection{Конечно-объёмная запись полной системы}

Интегрируем исходную систему по той же
пространственно-временной ячейке. Левая часть
преобразуется так же, как при выводе \eqref{eq:fv_exact}:
\[
V_i(\U_i^{n+1}-\U_i^n)
+\Delta t\left(
A_{i+1/2}\overline{\F}_{i+1/2}^{\,n}
-A_{i-1/2}\overline{\F}_{i-1/2}^{\,n}
\right)
=
\int_{t^n}^{t^{n+1}}
\int_{r_{i-1/2}}^{r_{i+1/2}}
\left(
A\Src-\bm K\,\frac{\partial(Au)}{\partial r}
\right)dr\,dt.
\]
Таким образом, потоковая часть сохраняет прежний вид.
Заменяя средние физические потоки численными,
получаем
\begin{equation}
\U_i^{n+1}
=
\U_i^n
-\frac{\Delta t}{V_i}
\left(
A_{i+1/2}\widehat{\F}_{i+1/2}^{\,n}
-A_{i-1/2}\widehat{\F}_{i-1/2}^{\,n}
\right)
+\Delta t\,\overline{\Src}_i^{\,n}
-\Delta t\,\overline{\bm C}_i^{\,n},
\label{eq:fv_full}
\end{equation}
где пока неизвестны интегральные средние
\[
\overline{\Src}_i^{\,n}
=
\frac{1}{\Delta t\,V_i}
\int_{t^n}^{t^{n+1}}
\int_{r_{i-1/2}}^{r_{i+1/2}}
A\Src\,dr\,dt,
\qquad
\overline{\bm C}_i^{\,n}
=
\frac{1}{\Delta t\,V_i}
\int_{t^n}^{t^{n+1}}
\int_{r_{i-1/2}}^{r_{i+1/2}}
\bm K\,\frac{\partial(Au)}{\partial r}\,dr\,dt.
\]
В отличие от потока, последний член нельзя
сразу представить разностью значений $\bm K Au$
на гранях: $\bm K$ зависит от решения.
Кроме того, при разрывах произведение
$\bm K\,\partial_r(Au)$ требует отдельного
определения.

\begin{figure}[h!]
    \centering
    \includegraphics[width=0.95\linewidth]{pics/grid_rt_correct.png}
    \caption{Схема сетки.}
\end{figure}

\subsection{Член расширения и сжатия: аппроксимация первого порядка}

Средний вклад неконсервативного члена в ячейке равен
\[
\overline{\bm C}_i^{\,n}
=
\frac{1}{\Delta t\,V_i}
\int_{t^n}^{t^{n+1}}
\int_{r_{i-1/2}}^{r_{i+1/2}}
\bm K\,\frac{\partial(Au)}{\partial r}\,dr\,dt.
\]
Заморозим коэффициент на старом временном слое:
\[
\bm K(r,t)\approx\bm K_i^n
=
\begin{pmatrix}
0&0&0&w_{+,i}^n/2&w_{-,i}^n/2
\end{pmatrix}^{\mathsf T}.
\]
После вынесения коэффициента пространственный интеграл
вычисляется через значения на границах ячейки:
\[
\overline{\bm C}_i^{\,n}
\approx
\frac{\bm K_i^n}{\Delta t\,V_i}
\int_{t^n}^{t^{n+1}}
\left[
A_{i+1/2}u_{i+1/2}(t)
-A_{i-1/2}u_{i-1/2}(t)
\right]dt.
\]
Аппроксимируя интеграл по времени значениями
на старом слое, получаем
\begin{equation}
\overline{\bm C}_i^{\,n}
\approx \bm K_i^n D_i^n,
\qquad
D_i^n=
\frac{
A_{i+1/2}\widehat u_{i+1/2}^{\,n}
-A_{i-1/2}\widehat u_{i-1/2}^{\,n}
}{V_i}.
\label{eq:compression_first_order}
\end{equation}
Соответствующая добавка к обновлению волновых энергий:
\[
\left.\Delta w_{\pm,i}\right|_{\text{сж/расш}}
=-\frac{\Delta t}{2}\,w_{\pm,i}^nD_i^n.
\]
При $D_i^n>0$ этот вклад уменьшает энергию волн,
при $D_i^n<0$ --- увеличивает её.

Остаётся определить скорость на каждой грани.
Для первого порядка используем состояния
$\U_L=\U_i^n$, $\U_R=\U_{i+1}^n$.
Рассмотрим два способа.

\subsubsection{Арифметическое среднее.}
Простейший выбор:
\[
\widehat u_{i+1/2}=\frac{u_L+u_R}{2}.
\]
Он согласован с точным значением на гладком решении,
но не учитывает направление распространения волн.

\subsubsection{Скорость промежуточного состояния HLL.}
Используем те же состояния, потоки и оценки скоростей,
что и при вычислении численного потока.
При $s_L<0<s_R$ промежуточное состояние имеет вид
\[
\U_*=
\frac{
s_R\U_R-s_L\U_L-(\F_R-\F_L)
}{s_R-s_L}.
\]
Тогда определим
\[
\widehat u_{i+1/2}=
\begin{cases}
u_L, & s_L\geq0,\\[2pt]
(\rho u)_*/\rho_*, & s_L<0<s_R,\\[2pt]
u_R, & s_R\leq0.
\end{cases}
\]
Для потока Русанова тот же выбор получается
из симметричных оценок $s_L=-\alpha$, $s_R=\alpha$:
\[
\U_*=
\frac{\U_L+\U_R}{2}
-\frac{\F_R-\F_L}{2\alpha},
\qquad
\widehat u_{i+1/2}=\frac{(\rho u)_*}{\rho_*}.
\]
Здесь предполагается $\rho_*>0$.

Предпочтём второй способ: скорость на грани
определяется по тому же приближённому решению задачи
Римана, что и численный поток.
Это мотивированный выбор дискретизации;
сам по себе он не гарантирует правильных условий
перехода волн через ударный фронт.


\subsection{Переход ко второму порядку}

Для второго порядка по пространству восстановим
в каждой ячейке линейные профили примитивных переменных
с ограничением наклонов. На грани используем
восстановленные состояния
\[
\bm q_L=\bm q_i+\bm\sigma_i(r_{i+1/2}-r_i),
\qquad
\bm q_R=\bm q_{i+1}
+\bm\sigma_{i+1}(r_{i+1/2}-r_{i+1}),
\qquad
\bm q=(\rho,u,p,w_+,w_-)^{\mathsf T}.
\]
По ним вычисляем как численный поток,
так и скорость $\widehat u_{i+1/2}$ выбранным выше способом.

На гладком решении центральная квадратура произведения
даёт пространственную аппроксимацию второго порядка:
\[
\frac1{V_i}
\int_{r_{i-1/2}}^{r_{i+1/2}}
\bm K\,\frac{\partial(Au)}{\partial r}\,dr
=
\bm K_i
\frac{
A_{i+1/2}\widehat u_{i+1/2}
-A_{i-1/2}\widehat u_{i-1/2}
}{V_i}
+O(h^2).
\]
Здесь предполагаются гладкая сетка и достаточно точное
вычисление геометрических величин.
Таким образом, форма \eqref{eq:compression_first_order}
сохраняется, но скорости на гранях вычисляются
по линейно восстановленным состояниям.

Для второго порядка по времени включим сжатие
в общий пространственный оператор:
\[
L_i(\U)=
-\frac{
A_{i+1/2}\widehat\F_{i+1/2}
-A_{i-1/2}\widehat\F_{i-1/2}
}{V_i}
+\Src_i-\bm K_iD_i.
\]
Затем применим двухстадийную схему Рунге--Кутты:
\[
\U_i^{(1)}=\U_i^n+\Delta t\,L_i(\U^n),
\]
\[
\U_i^{n+1}
=\frac12\U_i^n
+\frac12\left[
\U_i^{(1)}+\Delta t\,L_i(\U^{(1)})
\right].
\]
На каждой стадии заново выполняются реконструкция,
вычисление потоков, скоростей на гранях,
коэффициентов $\bm K_i$ и величин $D_i$.
Поэтому сжатие учитывается с той же временной точностью,
что и перенос.

Заявленный второй порядок относится к гладким решениям.
Вблизи разрывов ограничение наклонов снижает порядок,
а выбор смысла неконсервативного произведения
остаётся отдельным вопросом.



\subsection{Выбор скоростей для потоков HLL и LF}

Для вычисления численного потока нужны оценки
самой левой и самой правой скоростей распространения
возмущений через грань ячейки.
Обозначим их через $s_L$ и $s_R$.

Рассмотрим два способа их вычисления:
по характеристическим скоростям системы
и по приближённой формуле.

На грани $r_{i+1/2}$ имеются два состояния:
\[
\U_L=\U_i,\qquad \U_R=\U_{i+1}.
\]
Для первого порядка это средние значения в соседних
ячейках, для второго --- восстановленные значения
на грани на соответствующей стадии RK2.

Все необходимые величины вычисляются отдельно
для каждого из двух состояний:
\[
c_s^2=\frac{\gamma p}{\rho},
\qquad
V_A=\frac{b}{A_{i+1/2}\sqrt{\rho}}.
\]

\subsubsection{По характеристическим скоростям}

Рассмотрим характеристики подсистемы переноса
и сжатия. Градиентную часть отражения здесь
предполагаем обрабатываемой отдельно.

Одна характеристическая скорость равна $u$,
а остальные четыре имеют вид
\[
\lambda=u+x,
\]
где $x$ удовлетворяет уравнению
\begin{equation}
x^4-C_2x^2-C_1x+C_0=0,
\label{eq:speed_quartic}
\end{equation}
с коэффициентами
\[
C_2=c_s^2+V_A^2+\frac{3(w_++w_-)}{4\rho},
\]
\[
C_1=\frac{V_A(w_+-w_-)}{2\rho},
\qquad
C_0=V_A^2
\left(c_s^2+\frac{w_++w_-}{4\rho}\right).
\]

Для численного нахождения четырёх корней
$x_1,\ldots,x_4$ можно использовать
стандартный алгоритм поиска корней многочлена.
Например, они являются собственными значениями
матрицы
\[
\bm H=
\begin{pmatrix}
0&0&0&-C_0\\
1&0&0&C_1\\
0&1&0&C_2\\
0&0&1&0
\end{pmatrix}.
\]

Таким образом, для каждого состояния на грани
получаем пять характеристических скоростей:
\[
u,\quad u+x_1,\quad u+x_2,\quad
u+x_3,\quad u+x_4.
\]

Выберем минимальную и максимальную:
\[
\lambda_{\min}
=u+\min(0,x_1,x_2,x_3,x_4),
\]
\[
\lambda_{\max}
=u+\max(0,x_1,x_2,x_3,x_4).
\]

Дополнительно определим наибольшее отклонение
характеристической скорости от скорости газа:
\[
c_*=\max_{1\le k\le4}|x_k|.
\]
Эта величина используется только для задания
запаса при выборе границ скоростей.

Поскольку расчёт выполняется для обоих состояний,
получаем
\[
\lambda_{\min,L},\quad\lambda_{\max,L},\quad c_{*,L},
\]
\[
\lambda_{\min,R},\quad\lambda_{\max,R},\quad c_{*,R}.
\]

Теперь непосредственно зададим скорости
для численного потока:
\begin{equation}
\boxed{
\begin{aligned}
s_L&=\min\left(
\lambda_{\min,L}-\delta c_{*,L},\,
\lambda_{\min,R}-\delta c_{*,R}
\right),\\[3pt]
s_R&=\max\left(
\lambda_{\max,L}+\delta c_{*,L},\,
\lambda_{\max,R}+\delta c_{*,R}
\right).
\end{aligned}
}
\label{eq:speeds_exact}
\end{equation}

Здесь $\delta\ge0$ --- задаваемый относительный
запас. Например, для первоначальных расчётов
можно выбрать $\delta=0.1$.

Таким образом, $s_L$ оценивает скорость самого
левого возмущения, а $s_R$ --- самого правого.
Запас расширяет найденный интервал скоростей,
но не заменяет вычисление характеристик.

Для рассматриваемой подсистемы при
$\rho>0$, $p>0$, $w_\pm>0$ все пять
характеристических скоростей действительны.
Существенные мнимые части вычисленных корней
в этом случае свидетельствуют о численной
проблеме, которую необходимо исследовать.

\subsubsection{Приближённая оценка}

Вместо решения уравнения четвёртой степени
можно воспользоваться приближённой оценкой
максимальной скорости возмущений.

Обоснуем её с помощью двух предельных случаев.

\textbf{Первый случай: волновая энергия равна нулю.}

При $w_+=w_-=0$ уравнение
\eqref{eq:speed_quartic} принимает вид
\[
(x^2-c_s^2)(x^2-V_A^2)=0.
\]
Следовательно, характеристики равны
\[
u,\qquad u\pm c_s,\qquad u\pm V_A.
\]
Максимальное отклонение от скорости газа:
\[
c_{\max}=\max(c_s,V_A).
\]

\textbf{Второй случай: $V_A=0$.}

В этом формальном пределе характеристический
многочлен принимает вид
\[
x^2\left(
x^2-c_s^2-\frac{3(w_++w_-)}{4\rho}
\right)=0.
\]
Поэтому ненулевые относительно газа скорости
распространения равны
\[
x=\pm\sqrt{
c_s^2+\frac{3(w_++w_-)}{4\rho}
}.
\]

Таким образом, вклад волнового давления
в квадрат акустической скорости составляет
\[
c_w^2=\frac{3(w_++w_-)}{4\rho}.
\]

Объединим результаты двух предельных случаев.
В качестве приближённой максимальной скорости
относительно газа примем
\begin{equation*}
{
c_{\mathrm{fast}}=
\max\left(
V_A,\sqrt{c_s^2+\frac{3(w_++w_-)}{4\rho}}
\right).
}
\end{equation*}

Эта формула даёт правильные крайние скорости
в обоих рассмотренных предельных случаях.
Однако при одновременно ненулевых $V_A$
и $w_\pm$ она уже не является точной:
связь газовых и волновых возмущений
содержит дополнительные слагаемые.

В частности, она может занижать максимальную
характеристическую скорость.

Для каждого состояния на грани вычислим
$c_{\mathrm{fast},L}$ и $c_{\mathrm{fast},R}$.
Тогда непосредственно зададим
\begin{equation*}
{
\begin{aligned}
s_L&=\min\left(
u_L-c_{\mathrm{fast},L},\,
u_R-c_{\mathrm{fast},R}
\right),\\[3pt]
s_R&=\max\left(
u_L+c_{\mathrm{fast},L},\,
u_R+c_{\mathrm{fast},R}
\right).
\end{aligned}
}
\end{equation*}

В этих формулах дополнительный запас
не используется. При необходимости
можно заменить
\[
c_{\mathrm{fast}}
\longrightarrow
(1+\delta)c_{\mathrm{fast}}.
\]
Однако фиксированный процентный запас
не гарантирует охват всех характеристик
для произвольных допустимых состояний.

\subsubsection{Использование в HLL и LF}

В обоих рассмотренных вариантах результатом
являются две скорости на грани:
\[
\boxed{s_L,\qquad s_R.}
\]

Именно они используются в численном потоке HLL:
\[
\widehat{\F}_{\mathrm{HLL}}=
\begin{cases}
\F_L,&s_L\ge0,\\[3pt]
\displaystyle
\frac{
s_R\F_L-s_L\F_R+s_Ls_R(\U_R-\U_L)
}{s_R-s_L},
&s_L<0<s_R,\\[8pt]
\F_R,&s_R\le0.
\end{cases}
\]

Для потока Лакса--Фридрихса (Русанова)
достаточно одной оценки:
\[
\boxed{\alpha=\max(|s_L|,|s_R|),}
\]
после чего
\[
\widehat{\F}_{\mathrm{LF}}
=\frac{\F_L+\F_R}{2}
-\frac{\alpha}{2}(\U_R-\U_L).
\]

Те же оценки учитываются при выборе
шага по времени и пересчитываются
на каждой стадии RK2.

Основным вариантом выберем вычисление
характеристических скоростей, а приближённую
формулу оставим для сравнения.

Отметим два ограничения.

Во-первых, оценки характеристик только
левого и правого состояний не гарантируют,
что ими ограничены все скорости нелинейной
задачи Римана.

Во-вторых, характеристики были получены
с учётом неконсервативного члена сжатия,
тогда как поток HLL аппроксимирует
дивергентную часть системы. Поэтому
устойчивость их совместной дискретизации
необходимо проверять отдельно.

Наконец, если градиентную часть отражения
включать в главную матрицу полной системы,
характеристический многочлен изменяется.
В таком случае следует использовать
результаты полного характеристического анализа,
включая проверку гиперболичности.

\subsection{Выбор шага по времени: условие CFL}

Для каждой ячейки определим максимальный модуль
оценённых скоростей на обеих её гранях:
\[
a_i=\max\left(
|s_{L,i-1/2}|,\ |s_{R,i-1/2}|,\ 
|s_{L,i+1/2}|,\ |s_{R,i+1/2}|
\right).
\]
Здесь $s_L$ и $s_R$ --- скорости, выбранные выше
для вычисления численного потока. Шаг по времени зададим формулой
\[
\Delta t_{\mathrm{CFL}}
=C_{\mathrm{CFL}}\min_i\frac{\Delta r_i}{a_i},
\qquad
\Delta r_i=r_{i+1/2}-r_{i-1/2}.
\]
Для первоначальных расчётов примем
$C_{\mathrm{CFL}}=0.45$.
Тогда оценённое перемещение возмущения за шаг
не превышает $0.45$ полной ширины ячейки.
Запас, введённый при выборе $s_L,s_R$,
уже учитывается в величине $a_i$.

Для RK2 скорости пересчитываются на промежуточной
стадии. Если выбранный шаг нарушает ту же оценку,
расчёт шага повторяется с меньшим $\Delta t$.

Это ограничение относится к переносу.
Учёт источников и неконсервативного члена
может потребовать дополнительного уменьшения шага.

\subsection{Контроль положительности и шага по источникам}

Условие CFL ограничивает перенос между ячейками.
Источники могут уменьшать энергию непосредственно
внутри ячейки, поэтому для их явного обновления
может потребоваться меньший шаг.

Обозначим правые части, отвечающие за отражение
и диссипацию:
\[
S_+=-\Rref\sqrt{w_+w_-}-\Gamma_+w_+,
\qquad
S_-=\Rref\sqrt{w_+w_-}-\Gamma_-w_-.
\]
При явном обновлении
\[
w_\pm^{\,\mathrm{new}}=w_\pm+\Delta t\,S_\pm.
\]
Если $S_\pm<0$, источник расходует энергию волны.
Чтобы за один шаг он не израсходовал больше
имеющегося запаса, необходимо
\[
\Delta t\leq\frac{w_\pm}{|S_\pm|}.
\]
Поэтому для положительных $w_\pm$ примем
\[
\Delta t_{\mathrm{src}}
=
C_{\mathrm{src}}
\min_{\substack{i,\;\pm\\S_{\pm,i}<0}}
\frac{w_{\pm,i}}{|S_{\pm,i}|},
\qquad C_{\mathrm{src}}=0.5.
\]
Такой выбор позволяет источнику израсходовать
не более половины текущей энергии за подшаг.
Если убывающих источников нет, это ограничение
отсутствует. Общий шаг выбирается как
\[
\Delta t=\min(\Delta t_{\mathrm{CFL}},
             \Delta t_{\mathrm{src}}).
\]

После каждой стадии RK2 и отдельного обновления
источников проверяются
\[
\rho_i>0,\qquad w_{+,i}\geq0,\qquad w_{-,i}\geq0,
\]
\[
E_i-\frac{(\rho u)_i^2}{2\rho_i}
-w_{+,i}-w_{-,i}>0.
\]
Последнее неравенство означает положительность
тепловой энергии газа и давления.
Если хотя бы одно условие нарушено, весь шаг
повторяется с исходного состояния
с уменьшенным вдвое $\Delta t$.
Источники и ограничение на шаг пересчитываются
на каждой стадии.

Оценка $\Delta t_{\mathrm{src}}$ относится
к явному обновлению источников и не гарантирует
положительности всего совместного обновления.
При полунеявном интегрировании ограничение
определяется выбранной формулой источников.
Если уже при $w_\pm=0$ схема даёт $S_\pm<0$,
уменьшение шага не поможет: необходимо исправить
само обновление источника.

\subsection{Интегрирование источников по ячейке и по времени}

Вектор источников имеет вид
\[
\Src=
\begin{pmatrix}
0\\[3pt]
(p+P_A)\dfrac{A'}{A}-\dfrac{\rho}{r^2}\\[7pt]
-\dfrac{\rho u}{r^2}\\[7pt]
-\Rref\sqrt{w_+w_-}-\Gamma_+w_+\\[3pt]
+\Rref\sqrt{w_+w_-}-\Gamma_-w_-
\end{pmatrix}.
\]

Рассмотрим одну ячейку $I_i=[r_{i-1/2},r_{i+1/2}]$
объёма $V_i=\int_{I_i}A(r)\,dr$.
После интегрирования уравнений вклад источников за шаг равен
\[
\left.\Delta\U_i\right|_{\mathrm{ист}}
=
\int_{t^n}^{t^{n+1}}
\underbrace{\left[
\frac{1}{V_i}\int_{I_i}A(r)\Src(r,t)\,dr
\right]}_{\displaystyle \Src_i(t)}
dt.
\]
Сначала вычислим внутренний интеграл по каждой строке
при фиксированном $t=t^n$.

\subsubsection{Первая строка. Масса.}

\[
S_{\rho,i}(t^n)=\frac{1}{V_i}\int_{I_i}A(r)\cdot 0\,dr=0.
\]

\subsubsection{Вторая строка. Импульс.}

Умножим вторую строку на $A(r)$ и внесём под интеграл:
\[
S_{\rho u,i}(t^n)
=
\frac{1}{V_i}\int_{I_i}(p+P_A)A'(r)\,dr
-
\frac{1}{V_i}\int_{I_i}\rho\,\frac{A(r)}{r^2}\,dr.
\]

На слое $t^n$ примем газовое и волновое давления
постоянными внутри ячейки $I_i$ и равными
$p_i^n$ и $P_{A,i}^n$ соответственно. Тогда
\[
\frac{1}{V_i}\int_{I_i}(p+P_A)A'(r)\,dr
\;\approx\;
\frac{p_i^n+P_{A,i}^n}{V_i}\int_{I_i}A'(r)\,dr
=
(p_i^n+P_{A,i}^n)
\frac{A_{i+1/2}-A_{i-1/2}}{V_i}.
\]

В гравитационном члене на слое $t^n$ примем плотность
постоянной внутри ячейки $I_i$ и равной $\rho_i^n$.
Тогда её можно вынести за знак интеграла.
Складывая вклад давления и гравитации, получаем
\[
{
S_{\rho u,i}(t^n)\approx
\frac{1}{V_i}
\left[
(p_i^n+P_{A,i}^n)(A_{i+1/2}-A_{i-1/2})
-\rho_i^n
\int_{r_{i-1/2}}^{r_{i+1/2}}\frac{A(r)}{r^2}\,dr
\right].
}
\]
Оставшийся интеграл зависит только от заданной геометрии
и может быть вычислен заранее.

При $A(r)=r^2$ подынтегральная функция равна единице,
а объём ячейки равен
$V_i=(r_{i+1/2}^3-r_{i-1/2}^3)/3$. Поэтому
\[
S_{\rho u,i}(t^n)\approx
3\,
\frac{
(p_i^n+P_{A,i}^n)
\bigl(r_{i+1/2}^2-r_{i-1/2}^2\bigr)
-\rho_i^n\bigl(r_{i+1/2}-r_{i-1/2}\bigr)
}{
r_{i+1/2}^3-r_{i-1/2}^3
}
=
3\,
\frac{
(p_i^n+P_{A,i}^n)(r_{i+1/2}+r_{i-1/2})-\rho_i^n
}{
r_{i+1/2}^2+r_{i+1/2}r_{i-1/2}+r_{i-1/2}^2
}.
\]

\subsubsection{Третья строка. Полная энергия.}

На слое $t^n$ примем плотность импульса $\rho u$
постоянной внутри ячейки $I_i$ и равной $(\rho u)_i^n$.
Тогда
\[
S_{E,i}(t^n)
=
-\frac{1}{V_i}\int_{I_i}(\rho u)\frac{A(r)}{r^2}\,dr
\;\approx\;
-\frac{(\rho u)_i^n}{V_i}
\int_{r_{i-1/2}}^{r_{i+1/2}}\frac{A(r)}{r^2}\,dr.
\]
Оставшийся интеграл зависит только от заданной геометрии
и может быть вычислен заранее.

При $A(r)=r^2$ получаем
\[
S_{E,i}(t^n)\approx
-3(\rho u)_i^n
\frac{r_{i+1/2}-r_{i-1/2}}
     {r_{i+1/2}^3-r_{i-1/2}^3}
=
-\frac{3(\rho u)_i^n}
{r_{i+1/2}^2+r_{i+1/2}r_{i-1/2}+r_{i-1/2}^2}.
\]

\subsubsection{Четвёртая и пятая строки. Энергии волн.}

Подставим в обе строки
$\Gamma_\pm=d\sqrt{w_\mp/(A\rho)}$ и внесём множитель $A(r)$
под интегралы. Пока все преобразования точные:
\begin{align*}
S_{+,i}(t^n)
={}&-\frac{1}{V_i}\int_{I_i}
 A\,\Rref\sqrt{w_+w_-}\,dr
-\frac{d}{V_i}\int_{I_i}
 \sqrt{A}\,w_+\sqrt{\frac{w_-}{\rho}}\,dr,
\\[3pt]
S_{-,i}(t^n)
={}&+\frac{1}{V_i}\int_{I_i}
 A\,\Rref\sqrt{w_+w_-}\,dr
-\frac{d}{V_i}\int_{I_i}
 \sqrt{A}\,w_-\sqrt{\frac{w_+}{\rho}}\,dr.
\end{align*}
Уже здесь видно, что интеграл отражения один и тот же
в обеих строках. При их сложении он сокращается точно.

\emph{Диссипация.}
Известны средние $w_{+,i}^n$, $w_{-,i}^n$, $\rho_i^n$.
Чтобы вычислить интеграл, примем эти величины постоянными
внутри ячейки. Тогда все зависящие от решения множители
выносятся за знак интеграла:
\[
\frac{d}{V_i}\int_{I_i}
\sqrt{A(r)}\,w_\pm\sqrt{\frac{w_\mp}{\rho}}\,dr
\;\approx\;
\underbrace{\frac{d}{V_i}\int_{I_i}\sqrt{A(r)}\,dr}
_{\text{вычисляется по заданной геометрии}}
\,w_{\pm,i}^n\sqrt{\frac{w_{\mp,i}^n}{\rho_i^n}}
\;=:\;D_{\pm,i}^n.
\]

При $A(r)=r^2$ и $r>0$ имеем $\sqrt{A(r)}=r$, поэтому
\[
D_{\pm,i}^n
=
\frac{d}{V_i}
\frac{r_{i+1/2}^2-r_{i-1/2}^2}{2}
\,w_{\pm,i}^n\sqrt{\frac{w_{\mp,i}^n}{\rho_i^n}}
=
\frac{3d(r_{i+1/2}+r_{i-1/2})}
{2\bigl(r_{i+1/2}^2+r_{i+1/2}r_{i-1/2}+r_{i-1/2}^2\bigr)}
\,w_{\pm,i}^n\sqrt{\frac{w_{\mp,i}^n}{\rho_i^n}}.
\]

\emph{Отражение.}
В коэффициент $\Rref$ входит производная альфвеновской
скорости. Для её вычисления используем соседние ячейки
на одном временном слое:
\[
V_{A,j}^n=\frac{b}{A(r_j)\sqrt{\rho_j^n}},
\qquad j=i-1,\ i,\ i+1.
\]
На равномерной сетке примем
\[
\left(\frac{\partial V_A}{\partial r}\right)_i^n
\approx
\frac{V_{A,i+1}^n-V_{A,i-1}^n}{r_{i+1}-r_{i-1}}.
\]
Подставляя эту разность в замыкание, получаем
\[
\Rref_i^n=
\mathcal C(w_{+,i}^n,w_{-,i}^n)
\min\!\left\{
\left|\frac{V_{A,i+1}^n-V_{A,i-1}^n}{r_{i+1}-r_{i-1}}\right|,
\max(\Gamma_{+,i}^n,\Gamma_{-,i}^n)
\right\},
\qquad
\Gamma_{\pm,i}^n
=d\sqrt{\frac{w_{\mp,i}^n}{A(r_i)\rho_i^n}}.
\]

Примем $\Rref$, $w_+$ и $w_-$ постоянными внутри ячейки
и равными найденным ячеечным значениям. Тогда
\[
\frac{1}{V_i}\int_{I_i}
A(r)\Rref\sqrt{w_+w_-}\,dr
\;\approx\;
\Rref_i^n\sqrt{w_{+,i}^n w_{-,i}^n}
\frac{1}{V_i}\int_{I_i}A(r)\,dr
=
\Rref_i^n\sqrt{w_{+,i}^n w_{-,i}^n}
=:J_i^n.
\]
Последнее равенство следует из определения объёма $V_i$.

При $A(r)=r^2$ и $r>0$ используемые коэффициенты принимают вид
\[
V_{A,j}^n=\frac{b}{r_j^2\sqrt{\rho_j^n}},
\qquad
\Gamma_{\pm,i}^n
=\frac{d}{r_i}\sqrt{\frac{w_{\mp,i}^n}{\rho_i^n}},
\]
поэтому
\[
\Rref_i^n=
\mathcal C(w_{+,i}^n,w_{-,i}^n)
\min\!\left\{
\frac{|b|}{r_{i+1}-r_{i-1}}
\left|
\frac{1}{r_{i+1}^2\sqrt{\rho_{i+1}^n}}
-
\frac{1}{r_{i-1}^2\sqrt{\rho_{i-1}^n}}
\right|,
\frac{d}{r_i}
\sqrt{\frac{\max(w_{+,i}^n,w_{-,i}^n)}{\rho_i^n}}
\right\}.
\]
Вклад отражения по-прежнему равен
$J_i^n=\Rref_i^n\sqrt{w_{+,i}^n w_{-,i}^n}$.

С учётом диссипации четвёртая и пятая строки дают
\[
S_{+,i}(t^n)\approx-J_i^n-D_{+,i}^n,
\qquad
S_{-,i}(t^n)\approx+J_i^n-D_{-,i}^n.
\]

\subsubsection{Интеграл по времени.}

Обозначим через $\Src_i^n$ вектор из пяти только что
вычисленных строк. До сих пор использовались исключительно
данные слоя $t^n$. Для явного шага первого порядка остаётся
взять левый прямоугольник:
\[
\left.\Delta\U_i\right|_{\mathrm{ист}}
=\int_{t^n}^{t^{n+1}}\Src_i(t)\,dt
\;\approx\;\Delta t\,\Src_i^n.
\]
Именно на этом переходе источник считается постоянным
\emph{по времени} в течение шага.

\subsubsection{Как получить второй порядок.}

Для второго порядка по времени используем RK2 в форме
предиктора и корректора. Сначала по \emph{полной} правой части
сделаем предварительный шаг от $t^n$ до $t^{n+1}$.
Он даст состояния $\U_j^{(1)}$ во всех ячейках.
По ним повторим тот же расчёт пяти строк и получим
$\Src_i^{(1)}$. В частности, для отражения обе соседние
скорости $V_{A,i-1}^{(1)}$ и $V_{A,i+1}^{(1)}$ берутся
уже из предварительного решения. Тогда правило трапеций даёт
\[
\boxed{
\left.\Delta\U_i\right|_{\mathrm{ист}}
\approx
\frac{\Delta t}{2}\bigl(\Src_i^n+\Src_i^{(1)}\bigr).
}
\]
В этой форме RK2 предварительное состояние относится к концу
шага $t^{n+1}$. Окончательное состояние на том же слое
получается после усреднения правых частей.

\begin{figure}[htbp]
\centering
\begin{tikzpicture}[x=1.65cm,y=1cm,>=Stealth,font=\small]
  % Границы трёх пространственных ячеек.
  \foreach \x in {-0.5,0.5,1.5,2.5}
    \draw[gray!55,dashed] (\x,-0.35) -- (\x,2.45);

  % Нижняя строка -- известное решение,
  % верхняя -- предварительное решение в конце шага.
  \draw[gray] (-0.5,0) -- (2.5,0);
  \draw[gray] (-0.5,2) -- (2.5,2);
  \foreach \x/\j in {0/{i-1},1/{i},2/{i+1}} {
    \fill (\x,0) circle (1.6pt);
    \fill (\x,2) circle (1.6pt);
    \node[below] at (\x,-0.04) {$\U_{\j}^{n}$};
    \node[above] at (\x,2.04) {$\U_{\j}^{(1)}$};
  }

  \node[left] at (-0.55,0) {$t^n$};
  \node[left] at (-0.55,2) {$t^{n+1}$};
  \node[right] at (2.6,0) {$\longrightarrow\ \Src_i^n$};
  \node[right] at (2.6,2) {$\longrightarrow\ \Src_i^{(1)}$};

  % Скобка показывает шаблон производной на одном слое.
  \draw[<->] (0,0.55) -- (2,0.55)
    node[midway,fill=white,inner sep=2pt] {$2\Delta r$};
  \node at (1,1.15) {предварительный шаг RK2};
  \draw[->] (1,1.42) -- (1,1.84);
\end{tikzpicture}
\caption{В каждой строке рисунка источник вычисляется для
ячейки $i$. Значения слева и справа нужны для производной
$V_A$ в отражении. Для шага Эйлера используется нижняя строка,
для RK2 усредняются источники, полученные на обеих строках.}
\end{figure}

Для второго порядка \emph{по пространству} существенно,
насколько точно проинтегрирована каждая строка. На равномерной
сетке для гладкого решения среднее по ячейке отличается от
значения в её середине на $O(\Delta r^2)$. Это следует из
разложения около середины: интеграл линейного слагаемого
по симметричному интервалу равен нулю. Гладкий вес $A(r)$
сохраняет этот порядок. Поэтому на гладком участке замены,
сделанные в интегралах выше, уже дают второй порядок по
пространству. Центральная разность для $V_A$ имеет тот же
порядок. Для отражения этот вывод относится к участку,
где замыкание гладко. Чтобы второй порядок имела вся схема,
с тем же порядком необходимо аппроксимировать также потоки
на гранях и член расширения и сжатия.

\subsection{Возможная жёсткость источников и расщепление операторов}

В используемой схеме RK2 отражение и турбулентная диссипация
альфвеновских волн учитываются явно, как обычные источники.
Однако характерные времена этих процессов могут оказаться
существенно меньше временного шага, определяемого условием CFL.

Для диссипации характерные времена имеют вид
\[
\tau_{\mathrm{diss},\pm}=\frac{1}{\Gamma_\pm},
\qquad
\Gamma_\pm=d\sqrt{\frac{w_\mp}{A\rho}},
\]
где $d=2R_*/L_{\perp,\ph}$ --- введённый ранее
безразмерный параметр диссипации.

Поэтому необходимо контролировать величину
\[
\eta_{\mathrm{diss}}
=
\Delta t_{\mathrm{CFL}}
\max_i\left(\Gamma_{+,i},\Gamma_{-,i}\right).
\]
При $\eta_{\mathrm{diss}}\ll1$ явное интегрирование
диссипации не вызывает дополнительных затруднений.
Если же $\eta_{\mathrm{diss}}\gtrsim1$,
для сохранения устойчивости и положительности волновых
энергий может потребоваться уменьшение шага по времени.
Аналогичные ограничения могут возникать из-за быстрого
отражения волн, особенно в областях больших градиентов
альфвеновской скорости.

Предварительные оценки показывают, что для рассматриваемых
параметров диссипация, вероятно, не будет существенно
ограничивать шаг CFL. Поэтому первоначально все источники
предполагается интегрировать явно в рамках RK2.

Если в ходе расчётов обнаружится жёсткость источников,
можно перейти к расщеплению операторов. Например,
по схеме Стрэнга:
\[
\boxed{
\U^{n+1}
=
\mathcal S_{\Delta t/2}
\circ
\mathcal H_{\Delta t}
\circ
\mathcal S_{\Delta t/2}
(\U^n),
}
\]
где $\mathcal H_{\Delta t}$ обозначает шаг RK2
для пространственного оператора,
а $\mathcal S_{\Delta t/2}$ --- интегрирование
быстрых источников на полушаге.
Для устранения ограничения по времени диссипации
источники в этом случае следует интегрировать
неявно или полунеявно, причём с достаточной точностью,
чтобы сохранялся второй порядок всей схемы.


\section{Начальные условия}
\label{sec:initial_conditions}

\subsection{Безразмерная система и её коэффициенты}

Рассматриваем область $1\le r\le r_{\mathrm{out}}$.
Напомним безразмерные уравнения:
\[
\begin{aligned}
\frac{\partial(A\rho)}{\partial t}
+\frac{\partial(A\rho u)}{\partial r}
&=0,\\[4pt]
\frac{\partial(A\rho u)}{\partial t}
+\frac{\partial}{\partial r}
\left[A(\rho u^2+p+P_A)\right]
&=(p+P_A)A'-\frac{A\rho}{r^2},\\[4pt]
\frac{\partial(AE)}{\partial t}
+\frac{\partial}{\partial r}
\left[A\bigl(u(E+p+P_A)+V_A(w_+-w_-)\bigr)\right]
&=-\frac{A\rho u}{r^2},\\[4pt]
\frac{\partial(Aw_\pm)}{\partial t}
+\frac{\partial}{\partial r}
\left[A(u\pm V_A)w_\pm\right]
+\frac{w_\pm}{2}\frac{\partial(Au)}{\partial r}
&=\mp A\mathcal R\sqrt{w_+w_-}
-A\Gamma_\pm w_\pm.
\end{aligned}
\]
Здесь $A'=dA/dr$. Теплопроводность, радиационные потери
и дополнительный нагрев в этом варианте отсутствуют.

Связи между переменными:
\[
P_A=\frac{w_++w_-}{2},
\qquad
E=\frac{\rho u^2}{2}+\frac{p}{\gamma-1}+w_++w_-,
\]
\[
p=(\gamma-1)
\left(E-\frac{\rho u^2}{2}-w_+-w_-\right),
\qquad
T=\frac{p}{\theta\rho}.
\]

Используем масштабы
\[
R_*=R_\odot,\qquad
U_*=\sqrt{\frac{GM_\odot}{R_*}},\qquad
\rho_*=\rho_{\mathrm{ph}}^{\mathrm{phys}},\qquad
T_*=T_{\mathrm{ph}}^{\mathrm{phys}},\qquad
p_*=\rho_*U_*^2.
\]
Основные безразмерные коэффициенты равны
\[
\theta=\frac{k_{\mathrm B}T_*}{\mu m_uU_*^2},
\qquad
b=\frac{B_{\mathrm{ph}}^{\mathrm{phys}}}
{U_*\sqrt{\mu_0\rho_*}},
\qquad
d=\frac{2R_*}{L_{\perp,\mathrm{ph}}^{\mathrm{phys}}}.
\]
Коэффициент гравитации равен единице
благодаря выбору $U_*^2=GM_\odot/R_*$.

Альфвеновская скорость и частоты диссипации:
\[
V_A=\frac{b}{A\sqrt{\rho}},
\qquad
\Gamma_+=d\sqrt{\frac{w_-}{A\rho}},
\qquad
\Gamma_-=d\sqrt{\frac{w_+}{A\rho}}.
\]
Скорость передачи волновой энергии газу равна
\[
Q_w=\Gamma_+w_++\Gamma_-w_-.
\]
Она не добавляется в уравнение полной энергии:
диссипация перераспределяет энергию между волнами и газом.

Коэффициент отражения задаётся формулой
\[
\mathcal R
=
\mathcal C(w_+,w_-)
\min\!\left\{
\left|\frac{\partial V_A}{\partial r}\right|,
\max(\Gamma_+,\Gamma_-)
\right\},
\]
где
\[
\mathcal C(w_+,w_-)=
\begin{cases}
1-2\sqrt{w_-/w_+},
& w_+\ge4w_-,\quad w_+>0,\\[4pt]
2\sqrt{w_+/w_-}-1,
& w_-\ge4w_+,\quad w_->0,\\[4pt]
0,
& \text{в остальных случаях}.
\end{cases}
\]
В частности, $\mathcal C(0,0)=0$.
Производная альфвеновской скорости выражается через
геометрию и плотность:
\[
\frac{\partial V_A}{\partial r}
=
-V_A\left(
\frac{A'}{A}
+\frac{1}{2\rho}\frac{\partial\rho}{\partial r}
\right).
\]
Таким образом, $\mathcal C$, $\mathcal R$ и $\Gamma_\pm$
вычисляются по текущему решению и не являются
постоянными параметрами.

Для конкретного пробного расчёта возьмём основные
коэффициенты из безразмерного режима программы:
\[
\gamma=\frac53,\qquad
\mu=0.60,\qquad
\theta=0.0036,\qquad
b=1.50,\qquad
d=35,\qquad
\omega_0=0.08,\qquad
r_{\mathrm{out}}=30.
\]
Здесь
\[
\omega_0=
\frac{w_{+,\mathrm{ph}}^{\mathrm{phys}}(0)}{p_*}.
\]
Этот набор является вычислительным примером.
При выборе других размерных параметров фотосферы
коэффициенты следует пересчитать по приведённым формулам.

\subsection{Площадь магнитной трубки}

В общем случае зададим геометрию через положительную
функцию дополнительного расширения $f(r)$:
\[
A(r)=r^2\frac{f(r)}{f(1)},
\qquad f(r)>0.
\]
Тогда
\[
A(1)=1,\qquad
\frac{A'}{A}=\frac{2}{r}+\frac{f'}{f},
\]
\[
V_A=\frac{b\,f(1)}{r^2f(r)\sqrt{\rho}},
\qquad
\Gamma_\pm
=\frac{d}{r}
\sqrt{\frac{f(1)w_\mp}{f(r)\rho}}.
\]
Функция $f(r)$ задаётся заранее и не зависит от времени.

В рассматриваемом радиальном варианте выбираем
$f(r)\equiv1$. Поэтому
\[
A(r)=r^2,\qquad
A'(r)=2r,\qquad
\frac{A'}{A}=\frac2r,
\]
\[
V_A=\frac{b}{r^2\sqrt{\rho}},
\qquad
\Gamma_\pm=\frac{d}{r}\sqrt{\frac{w_\mp}{\rho}},
\]
\[
\mathcal R
=
\mathcal C(w_+,w_-)
\min\!\left\{
V_A\left|
\frac2r+\frac{1}{2\rho}\frac{\partial\rho}{\partial r}
\right|,
\frac{d}{r}
\sqrt{\frac{\max(w_+,w_-)}{\rho}}
\right\}.
\]

\subsection{Согласование с граничными условиями}

На фотосфере $r_{\mathrm{ph}}=1$ задаём
\[
\rho_{\mathrm{ph}}(t)=1,\qquad
T_{\mathrm{ph}}(t)=1,\qquad
p_{\mathrm{ph}}(t)=\theta,\qquad
w_{+,\mathrm{ph}}(t)=\omega(t).
\]
Для поиска стационарного течения используем
$\omega(t)\equiv\omega_0$.
Скорость $u_{\mathrm{ph}}(t)$ и энергию
$w_{-,\mathrm{ph}}(t)$ определяем по информации из области
с помощью выбранного граничного алгоритма.

На внешней границе свободный выход допустим,
если все характеристики направлены из области.
При наличии входящих характеристик
для них необходимо задать внешние данные.

\subsection{Начальная скорость и плотность}

Используем следующую форму начальных профилей:
\[
u(r,0)
=
u_{\mathrm{ph}}^0
+1.45\left(1-\exp\!\left[-\frac{r-1}{3.5}\right]\right),
\qquad
u_{\mathrm{ph}}^0=0.015.
\]
Для выбранных коэффициентов
\[
c_{s,\mathrm{ph}}^0=\sqrt{\gamma\theta}
=\sqrt{0.006},
\qquad
V_{A,\mathrm{ph}}^0=b=1.50.
\]
Поэтому начальная скорость на фотосфере
меньше обеих этих скоростей.

Плотность задаём формулой
\[
\rho(r,0)
=
\max\!\left\{
\frac{u_{\mathrm{ph}}^0}{A(r)u(r,0)}
+
\left(1-\frac{u_{\mathrm{ph}}^0}{u(r,0)}\right)
\exp\!\left[-\frac{r-1}{0.15}\right],
\ 10^{-8}
\right\}.
\]
На фотосфере получаем
\[
u(1,0)=u_{\mathrm{ph}}^0,\qquad
\rho(1,0)=1.
\]
Во внешней части области экспоненциальная добавка мала,
и до срабатывания нижнего ограничения плотности
\[
A(r)\rho(r,0)u(r,0)\simeq u_{\mathrm{ph}}^0.
\]
Это условие приблизительно постоянного массового расхода
для начального профиля, а не заранее заданный расход
будущего стационарного решения.

\subsection{Начальная температура и давление}

Зададим температурный профиль с максимумом:
\[
T(r,0)
=
1+25
\left(1-\exp\!\left[-\frac{r-1}{0.3}\right]\right)
r^{-1/4}.
\]
Он удовлетворяет условию $T(1,0)=1$,
возрастает над фотосферой,
достигает максимума и затем убывает.
Число $25$ задаёт амплитуду добавки,
а не само значение максимальной температуры.

Давление вычисляем по уравнению состояния:
\[
p(r,0)=\theta\rho(r,0)T(r,0).
\]
Следовательно,
\[
p(1,0)=\theta.
\]
Размерная температура равна
\[
T^{\mathrm{phys}}(r,0)=T_*T(r,0).
\]
Поэтому при изменении фотосферного масштаба $T_*$
размерная высота температурного максимума
также изменяется.

\subsection{Начальные энергии волн}

Выберем профили
\[
w_+(r,0)=\omega_0 A(r)^{-1.2},
\]
\[
w_-(r,0)
=
0.02\,w_+(r,0)
\left(1-\exp\!\left[-\frac{r-1}{0.5}\right]\right).
\]
В радиальной геометрии они принимают вид
\[
w_+(r,0)=\omega_0 r^{-2.4},
\]
\[
w_-(r,0)
=
0.02\,\omega_0 r^{-2.4}
\left(1-\exp\!\left[-\frac{r-1}{0.5}\right]\right).
\]
На фотосфере
\[
w_{+,\mathrm{ph}}^0=\omega_0,\qquad
w_{-,\mathrm{ph}}^0=0.
\]
При $r>1$ обе энергии положительны,
а отношение $w_-/w_+$ возрастает от нуля к $0.02$.
Нулевое значение встречной энергии на фотосфере
задаётся только в начальный момент;
далее оно определяется граничным алгоритмом.

\subsection{Полная энергия и заполнение ячеек}

Начальную полную энергию вычисляем как
\[
E(r,0)
=
\frac{\rho(r,0)u(r,0)^2}{2}
+\frac{p(r,0)}{\gamma-1}
+w_+(r,0)+w_-(r,0).
\]
Как в программе, значения для инициализации
берём в центрах ячеек:
\[
\bm U_i^0=
\begin{pmatrix}
\rho(r_i,0)\\
\rho(r_i,0)u(r_i,0)\\
E(r_i,0)\\
w_+(r_i,0)\\
w_-(r_i,0)
\end{pmatrix}.
\]
Это приближение объёмных средних
значениями в центрах ячеек.
При кусочно-постоянном представлении
внутри ячейки используется состояние $\bm U_i^0$.
Если хранятся переменные $A\bm U$,
начальный вектор равен $A(r_i)\bm U_i^0$.

Заданные профили описывают начальный ускоряющийся поток
с горячей областью над фотосферой.
Они не удовлетворяют всем стационарным балансам.
Сохранение горячей короны и установление течения
должны определяться последующей эволюцией
по полной системе уравнений.


\section{Растянутая пространственная сетка}
\label{sec:stretched_grid}

\subsection{Зачем растягивать сетку}

Вблизи фотосферы плотность, температура и волновые
энергии могут изменяться на малых пространственных
масштабах. Поэтому здесь нужны небольшие ячейки.
Во внешней части области для описания плавных профилей
можно использовать более крупные ячейки.

Выберем геометрически растущие размеры ячеек:
\[
\Delta r_{i+1}=q\Delta r_i,\qquad q=1.008.
\]
Каждая следующая ячейка шире предыдущей на $0.8\%$.
Такое растяжение позволяет сосредоточить разрешение
возле фотосферы, сохраняя умеренное число ячеек
во всей области.

Число $q=1.008$ является параметром сетки,
а не физической модели.
Достаточность выбранного разрешения
проверяется сравнением расчётов на разных сетках.

\subsection{Построение сетки}

Пусть область $[1,r_{\mathrm{out}}]$ разбита
на $N$ ячеек:
\[
I_i=[r_{i-1/2},r_{i+1/2}],
\qquad
\Delta r_i=r_{i+1/2}-r_{i-1/2}.
\]
Задаём
\[
\Delta r_i=\Delta r_1q^{i-1}.
\]
Сумма размеров ячеек должна совпадать
с длиной расчётной области:
\[
r_{\mathrm{out}}-1
=\sum_{i=1}^{N}\Delta r_i
=\Delta r_1\frac{q^N-1}{q-1}.
\]
Следовательно,
\[
\Delta r_1
=(r_{\mathrm{out}}-1)\frac{q-1}{q^N-1}.
\]
Координаты граней и центров:
\[
r_{i+1/2}
=1+\Delta r_1\frac{q^i-1}{q-1},
\qquad i=0,\ldots,N,
\]
\[
r_i=\frac{r_{i-1/2}+r_{i+1/2}}{2},
\qquad i=1,\ldots,N.
\]
При $q=1$ используется равномерная сетка:
\[
\Delta r_i=\frac{r_{\mathrm{out}}-1}{N}.
\]

Например, при
\[
N=600,\qquad r_{\mathrm{out}}=30,\qquad q=1.008
\]
получаем
\[
\Delta r_1\simeq0.001963,\qquad
\Delta r_N\simeq0.2321,\qquad
\frac{\Delta r_N}{\Delta r_1}\simeq118.3.
\]
Для равномерной сетки из тех же $600$ ячеек
\[
\Delta r_{\mathrm{uniform}}\simeq0.04833.
\]
Таким образом, возле фотосферы шаг примерно
в $24.6$ раза меньше, чем на равномерной сетке
с тем же числом ячеек.
Все приведённые расстояния безразмерны;
размерный шаг равен $R_*\Delta r_i$.

\subsection{Влияние на уравнения и численные потоки}

Растяжение сетки не меняет исходные
дифференциальные уравнения.
Мы продолжаем использовать ту же координату $r$,
но выбираем неравномерно расположенные грани ячеек.
Дополнительных физических слагаемых
или коэффициента $q$ в уравнениях не возникает.

Меняются геометрические величины в дискретной схеме.
Для записи
\[
\frac{\partial(A\bm U)}{\partial t}
+\frac{\partial(A\bm F)}{\partial r}
+\bm K\frac{\partial(Au)}{\partial r}
=A\bm S
\]
вводим объём и среднее состояние ячейки:
\[
V_i=\int_{r_{i-1/2}}^{r_{i+1/2}}A(r)\,dr,
\qquad
\bm U_i
=\frac1{V_i}\int_{I_i}A(r)\bm U(r)\,dr.
\]
При $A(r)=r^2$
\[
V_i=\frac{r_{i+1/2}^3-r_{i-1/2}^3}{3}.
\]

Например, явное обновление первого порядка имеет вид
\[
\begin{aligned}
\bm U_i^{n+1}
={}&\bm U_i^n
-\frac{\Delta t}{V_i}
\left(
A_{i+1/2}\widehat{\bm F}_{i+1/2}^n
-A_{i-1/2}\widehat{\bm F}_{i-1/2}^n
\right)
\\
&+\Delta t\,\bm S_i^n
-\frac{\Delta t}{V_i}\bm K_i^n
\left(
A_{i+1/2}\widehat u_{i+1/2}^n
-A_{i-1/2}\widehat u_{i-1/2}^n
\right),
\end{aligned}
\]
где
\[
A_{i\pm1/2}=A(r_{i\pm1/2}),
\qquad
\bm S_i^n\simeq
\frac1{V_i}\int_{I_i}A(r)\bm S(r,t^n)\,dr.
\]
Формула численного потока HLL или LF сохраняется.
В ней используются состояния слева и справа
от соответствующей грани.

Один и тот же поток через общую грань
входит в балансы соседних ячеек с противоположными
знаками. Поэтому неравномерность сетки
не нарушает консервативность потоковой части.

\subsection{Пространственные производные}

В разностных формулах необходимо использовать
фактические расстояния между центрами.
Например,
\[
\left(\frac{\partial V_A}{\partial r}\right)_i
\simeq
\frac{V_{A,i+1}-V_{A,i-1}}{r_{i+1}-r_{i-1}}.
\]
Однако на неравномерной сетке эта формула
в общем случае не имеет второго порядка
в точке $r_i$.

Если требуется трёхточечная аппроксимация
второго порядка для гладкого профиля,
обозначим расстояния до соседних центров:
\[
\delta_-=r_i-r_{i-1},
\qquad
\delta_+=r_{i+1}-r_i.
\]
Тогда
\[
\left(\frac{\partial V_A}{\partial r}\right)_i
\simeq
-\frac{\delta_+}{\delta_-(\delta_-+\delta_+)}V_{A,i-1}
+\frac{\delta_+-\delta_-}{\delta_-\delta_+}V_{A,i}
+\frac{\delta_-}{\delta_+(\delta_-+\delta_+)}V_{A,i+1}.
\]
Эта производная используется при вычислении
коэффициента отражения $\mathcal R$.
Аналогично при реконструкции профилей
наклоны вычисляются с учётом реальных расстояний
между центрами ячеек.

\subsection{Шаг по времени и точность}

Для переноса используем оценку шага вида
\[
\Delta t_{\mathrm{CFL}}
=
C_{\mathrm{CFL}}
\min_i\frac{\Delta r_i}{a_i},
\]
где $a_i$ --- оценка максимального модуля
характеристической скорости в окрестности ячейки,
а $C_{\mathrm{CFL}}$ выбирается с учётом
применяемой схемы.
Ограничения от источников и контроля положительности
учитываются дополнительно.

Мелкие ячейки возле фотосферы могут ограничивать
общий шаг по времени.
Поэтому растяжение не обязательно ускоряет расчёт
по сравнению с равномерной сеткой из тех же $N$ ячеек:
оно прежде всего перераспределяет разрешение.
Зато равномерная сетка с таким же малым шагом,
как $\Delta r_1$, потребовала бы в данном примере
около $14777$ ячеек вместо $600$.

Крупные внешние ячейки хуже разрешают короткие волны
и узкие структуры и обычно усиливают их численное
размазывание. Поэтому при исследовании нестационарных
возмущений нужно проверять разрешение
во всей области, а не только возле фотосферы.

При проверке сходимости следует уменьшать размеры
всех ячеек. Например, каждую ячейку исходной сетки
можно разделить пополам.
Одно увеличение $N$ при фиксированных $q>1$
и границах области не гарантирует измельчения
внешней части сетки.

При переходе на растянутую сетку необходимо:
\begin{itemize}
    \item Пересчитать координаты граней и центров,
    размеры $\Delta r_i$ и объёмы $V_i$ всех ячеек.

    \item Вычислять $A$, $A'$ и площади граней
    в новых точках сетки.

    \item В разностях для $\partial V_A/\partial r$
    учитывать реальные расстояния между центрами;
    для второго порядка использовать трёхточечную
    формулу на неравномерной сетке.

    \item При линейной реконструкции учитывать
    местные расстояния между центрами и гранями.
    Для кусочно-постоянных состояний это не требуется.

    \item В условии CFL использовать локальные шаги
    $\Delta r_i$, а при проверке оснований граничных
    характеристик --- ширину первой ячейки.

    \item В потоковом обновлении и интегрировании
    источников использовать новые границы
    и объёмы ячеек.
\end{itemize}
Исходные уравнения, физические коэффициенты,
формулы потоков HLL/LF и стадии RK2 не меняются.








\newpage
\appendix
\section{Явные формулы характеристических скоростей}
\label{app:characteristic_roots}

Рассмотрим полный характеристический полином
\[
\det(\lambda\mathsf I-\mathsf J)=zP(z),
\qquad z=\lambda-u,
\]
где
\[
P(z)
=
(z^2-c_s^2)(z^2-V_A^2)
-\frac{
3Wz^2+2V_A\Delta Wz-V_A^2(W+2h)
}{4\rho},
\qquad
W=w_++w_-,
\qquad
\Delta W=w_+-w_-.
\]
Предполагаем, что все пять корней полного характеристического
полинома вещественны, с учётом кратности.
Никаких предположений о малости волновой энергии
или коэффициента $h$ не вводится.

\subsection{Приведение к уравнению четвёртой степени}

Только в пределах настоящего приложения обозначим
коэффициенты полинома через
\[
\begin{aligned}
a&=c_s^2+V_A^2+\frac{3W}{4\rho},\\[3pt]
b&=\frac{V_A\Delta W}{2\rho},\\[3pt]
c&=V_A^2\left(c_s^2+\frac{W+2h}{4\rho}\right).
\end{aligned}
\]
Тогда характеристическое уравнение принимает вид
\[
z\left(z^4-az^2-bz+c\right)=0.
\]
Один корень известен сразу:
\[
z_0=0,\qquad \lambda_0=u.
\]
Остальные четыре корня определяются уравнением
\begin{equation}
z^4-az^2-bz+c=0.
\label{app:quartic_equation}
\end{equation}

\subsection{Общий случай: $b\ne0$}

Введём положительную величину $m$ и представим полином
как разность квадратов:
\[
\begin{aligned}
z^4-az^2-bz+c
={}&
\left(z^2+\frac{m-a}{2}\right)^2
-\left(\sqrt m\,z+\frac{b}{2\sqrt m}\right)^2
+
c-\frac{(m-a)^2}{4}+\frac{b^2}{4m}.
\end{aligned}
\]
Чтобы последняя строка обратилась в нуль, необходимо
\[
m(m-a)^2-4cm-b^2=0.
\]
Таким образом, $m$ должно удовлетворять
вспомогательному кубическому уравнению
\begin{equation}
m^3-2am^2+(a^2-4c)m-b^2=0.
\label{app:cubic_resolvent}
\end{equation}

Выбираем наибольший вещественный корень
уравнения~\eqref{app:cubic_resolvent}.
При четырёх вещественных корнях
уравнения~\eqref{app:quartic_equation} и $b\ne0$
этот корень положителен.

Для пояснения обозначим корни квартета через
$r_1,r_2,r_3,r_4$.
Поскольку кубический член отсутствует, по теореме Виета получаем
\[
r_1+r_2+r_3+r_4=0.
\]
Три корня кубического
уравнения~\eqref{app:cubic_resolvent} равны
\[
(r_1+r_2)^2,\qquad
(r_1+r_3)^2,\qquad
(r_1+r_4)^2.
\]
Следовательно, они вещественны и неотрицательны.
При $b\ne0$ их произведение равно $b^2>0$,
поэтому все три положительны.

Выбранный наибольший корень можно записать явно (тригонометрический метод решения кубического уравнения, или формула Виета для кубического уравнения):
\[
m=
\frac{2a}{3}
+\frac{2\sqrt{a^2+12c}}{3}
\cos\!\left[
\frac13\arccos\!\left(
\frac{-2a^3+72ac+27b^2}
{2(a^2+12c)^{3/2}}
\right)
\right].
\]
При рассматриваемых предположениях
$a^2+12c>0$, а аргумент $\arccos$ принадлежит $[-1,1]$.

После выбора $m$ получаем точное разложение
\[
\begin{aligned}
z^4-az^2-bz+c
={}&
\left(
z^2-\sqrt m\,z+\frac{m-a}{2}
-\frac{b}{2\sqrt m}
\right)\cdot
\left(
z^2+\sqrt m\,z+\frac{m-a}{2}
+\frac{b}{2\sqrt m}
\right).
\end{aligned}
\]
Решая два квадратных уравнения, находим
\begin{equation}
\begin{aligned}
z_{1,2}
&=
\frac{
\sqrt m\ \pm
\sqrt{\,2a-m+\dfrac{2b}{\sqrt m}\,}
}{2},
\\[6pt]
z_{3,4}
&=
\frac{
-\sqrt m\ \pm
\sqrt{\,2a-m-\dfrac{2b}{\sqrt m}\,}
}{2}.
\end{aligned}
\label{app:four_explicit_roots}
\end{equation}
В каждой строке берутся оба знака.
При вещественности всех четырёх корней
подкоренные выражения неотрицательны.
Нумерация в~\eqref{app:four_explicit_roots}
не задаёт порядка корней по величине.

\subsection{Биквадратный случай: $b=0$}

Если $b=0$, уравнение~\eqref{app:quartic_equation}
решается заменой $y=z^2$:
\[
y^2-ay+c=0.
\]
Поэтому
\[
\begin{aligned}
z_{1,2}
&=\pm\sqrt{\frac{a+\sqrt{a^2-4c}}{2}},\\[5pt]
z_{3,4}
&=\pm\sqrt{\frac{a-\sqrt{a^2-4c}}{2}}.
\end{aligned}
\]
Предположение о вещественности всех четырёх корней
означает
\[
a^2-4c\ge0,
\qquad
\frac{a+\sqrt{a^2-4c}}2\ge0,
\qquad
\frac{a-\sqrt{a^2-4c}}2\ge0.
\]
Коэффициент $h$ сохраняется в этих формулах через $c$.

\subsection{Все пять характеристических скоростей}

После вычисления корней квартета получаем
\[
\lambda_0=u,\qquad
\lambda_j=u+z_j,\qquad j=1,2,3,4.
\]
При необходимости скорости сортируются по возрастанию.

Формулы учитывают все слагаемые исходного полинома.
Совпадающие корни выдаются с учётом кратности.
Вещественность всех корней сама по себе не означает
строгой гиперболичности: для неё все пять скоростей
должны быть различными.

При численной реализации около кратных корней
возможна потеря точности вследствие вычитания близких
чисел. Поэтому найденные корни следует проверять
подстановкой в исходный полином.
\end{document}
